% !TeX root = ../../Book4.tex
% 纲要标题：## 第7章　线性辛几何、余切丛与相容结构
% 纲要位置：几何学计划纲要.md，第 3741 行。
% 正文状态：初稿；尚未正式审核。

% 纲要内容开始
% 
% > 本章目标与证明范围：在辛基础中建立相容概复结构的存在和可缩性，不要求先经过全纯曲线或 Kähler。闭性、非退化性和可积性始终是不同条件。
% 
% 纲要内容结束

\chapter{线性辛几何、余切丛与相容结构}
\label{g4:ch:07}
\BookChapterNumber{g4:ch:07}

\begin{chapterabstract}
本章从非退化反对称配对进入辛几何。辛基把所有同维线性辛空间化为同一模型，辛正交补刻画各类特殊子空间，并给出线性约化。余切丛上的典范一形式把线性模型推广到流形，其图像与闭一形式之间的对应说明闭性怎样进入几何。通过正定算子的极分解，构造与辛形式相容的几乎复结构，并给出相容结构空间的显式收缩。最后从体积、上同调及陈类得到整体约束，用具体计算区分非退化性、闭性与复结构的可积性。
\end{chapterabstract}

\begin{learninggoals}
\begin{itemize}
	\item 构造辛基，计算辛正交补，判定各向同性、余各向同性及 Lagrange 子空间，并完成线性约化。
	\item 证明辛群保持定向，写出其到酉群的形变收缩，理解线性辛群的基本拓扑。
	\item 从余切向量的自然配对定义典范一形式，核对局部坐标、余切提升与一形式图像的符号。
	\item 从任意辅助正定内积构造相容几乎复结构，证明构造的光滑性、自然性及相容结构空间的可缩性。
	\item 把点态构造推广到流形及光滑参数族，并说明相容结构的选择为何不影响辛切丛的陈类。
	\item 用体积排除闭辛流形上的恰当辛形式，计算基本例子，并识别相容但不可积的几乎复结构。
\end{itemize}
\end{learninggoals}

在平面上，有向面积给两个向量一个反对称配对。它既能检验两个方向是否独立，也能把一个方向变成作用在其他向量上的线性函数。高维辛几何保留的是这种非退化配对：每个非零向量总能找到与它配对不为零的向量。一个向量与自身的配对始终为零，因此辛形式不能直接测长度。长度会在选择相容结构之后出现。

\ProseParagraphBreak

把同样的配对放到流形的每根切空间上，还需要要求二形式闭。余切丛提供了最自然的例子：底点的一个切向量与该点的一份余切数据直接配对，产生典范一形式；再取负外微分，便得到辛形式。这个构造不需要先有度量或复坐标。本编可以从这里独立进入：所需基础是第一卷的向量丛、微分形式与积分，以及第二卷的基本群和示性类；前六章的复曲线理论提供比较背景。

\ProseParagraphBreak

相容几乎复结构把辛形式与正定度量联系起来。关键问题有两个：这样的结构能否光滑地选遍流形，以及两次选择之间是否存在不可消除的差别。本章用一个统一公式回答它们。先由辅助度量把辛形式写成斜自伴算子，再用正平方根除去其伸缩部分。正定度量可以线性插值，经过这个归一化后便得到相容结构空间的收缩。可缩性还保证切丛陈类的定义只依赖辛形式。

\ProseParagraphBreak

本章约定标准辛形式为 \(\sum_j\mathrm{d}q^j\wedge\mathrm{d}p_j\)，余切丛的典范辛形式为 \(-\mathrm{d}\theta\)，相容度量为 \(g(u,v)=\omega(u,Jv)\)。这些符号将在后续动力学中继续使用。线性模型、余切丛与相容结构可对照 Cannas da Silva 的正式综述第1.1—1.3、2.1及3.1—3.2节；本章在正文中给出所需构造和证明。\cite[第1.1--1.3、2.1、3.1--3.2节]{Cannas2006Symplectic}

% !TeX root = ../../Book4.tex
% 纲要标题：### 7.1 线性辛空间
% 纲要位置：几何学计划纲要.md，第 3745 行。
% 正文状态：初稿；尚未正式审核。

% 纲要内容开始
% 
% * 辛基与交替型标准形
% * 各向同性、余各向同性和 Lagrange
% * 辛约化的线性模型
% * 辛群及其基本拓扑
% 
% 纲要内容结束

\section{线性辛空间}
\label{g4:sec:0701}

在平面上，行列式 \(qp'-pq'\) 测量两个向量张成的平行四边形的有向面积。辛形式把这种配对推广到任意偶数维空间：它仍然交换两个变量就变号，而且每个非零向量都能与某个向量配出非零面积。本节从这一线性结构出发，说明标准坐标、特殊子空间和约化的含义，最后研究保持这种配对的线性变换群。流形上的辛几何将逐点使用这些结论。

\subsection{辛基与交替型标准形}
\label{g4:subsec:0701:01}

以下向量空间均为有限维实向量空间。双线性型 \(\omega:V\times V\to\mathbb R\) 称为交替型，是指 \(\omega(v,v)=0\) 对所有 \(v\) 成立；在实数域上，这等价于 \(\omega(u,v)=-\omega(v,u)\)。其根空间为
\[
	\operatorname{rad}\omega=\{v\in V:\omega(v,w)=0\text{ 对所有 }w\in V\}.
\]
这里的“根”记录了配对无法检测的方向。

\begin{definition}\label{g4:def:0701:symplectic-space}
	若交替型 \(\omega\) 的根空间为零，就称 \(\omega\) 非退化，并称 \((V,\omega)\) 为\term[xinxianxingkongjian]{辛线性空间}[symplectic vector space]，简称辛空间。等价地，线性映射
	\[
		\omega^\flat:V\longrightarrow V^*,\qquad
		\omega^\flat(v)(w)=\omega(v,w)
	\]
	是同构。
\end{definition}

内积也给出 \(V\to V^*\) 的同构，但内积满足 \(\langle v,v\rangle>0\)，而交替型总有 \(\omega(v,v)=0\)。因此，辛形式的非退化性要通过向量与其他方向的配对来判断。

\begin{theorem}[交替型的标准形]\label{g4:thm:0701:symplectic-basis}
	设 \(\omega\) 是 \(V\) 上的交替型。存在一组基
	\[
		e_1,\ldots,e_r,f_1,\ldots,f_r,z_1,\ldots,z_k
	\]
	使得
	\[
		\omega(e_i,e_j)=\omega(f_i,f_j)=0,\qquad
		\omega(e_i,f_j)=\delta_{ij},
	\]
	且 \(z_1,\ldots,z_k\) 张成 \(\operatorname{rad}\omega\)。特别地，交替型的秩为偶数 \(2r\)。若 \(\omega\) 非退化，则 \(\dim V=2n\) 为偶数，上述基中没有 \(z_j\)，并称 \((e_1,\ldots,e_n,f_1,\ldots,f_n)\) 为\term[xinji]{辛基}[symplectic basis]。
\end{theorem}
\begin{proof}
	对 \(\dim V\) 归纳。若 \(\omega=0\)，任一组基都是根空间的基，结论成立。否则可以选取 \(e,f\) 使 \(\omega(e,f)\ne0\)，将 \(f\) 乘以适当常数后有 \(\omega(e,f)=1\)。令 \(H=\operatorname{span}(e,f)\)，并定义
	\[
		P(v)=\omega(v,f)e-\omega(v,e)f.
	\]
	直接计算得 \(P(e)=e\)、\(P(f)=f\)，而 \(v-P(v)\) 与 \(e,f\) 的配对均为零。于是
	\[
		V=H\oplus H',\qquad
		H'=\{v:\omega(v,e)=\omega(v,f)=0\},
	\]
	且两个直和分量之间没有交叉配对。对 \(H'\) 上的限制型应用归纳假设，连同 \(e,f\) 就得到所需的基。

	在这组基中，每个 \(e_i,f_i\) 产生一个可逆的 \(2\times2\) 交替块；其余 \(z_j\) 与所有基向量的配对为零。故根空间恰由 \(z_j\) 张成，秩为 \(2r\)。非退化时 \(k=0\)，于是维数为 \(2r\)。
\end{proof}

给定辛基，以 \(q^i,p_i\) 表示对应的线性坐标。坐标函数的微分在每个点都可视为同一对偶向量，因而标准形写成
\begin{equation}\label{g4:eq:0701:standard-form}
	\omega_0=\sum_{i=1}^n\mathrm{d}q^i\wedge\mathrm{d}p_i,
	\qquad
	\omega_0(u,v)=u^{\mathsf T}\Omega v,
	\qquad
	\Omega=\begin{pmatrix}0&I_n\\-I_n&0\end{pmatrix}.
\end{equation}
向量的分量始终按全部 \(q\) 分量在前、全部 \(p\) 分量在后排列。两个辛空间之间的线性同构 \(F:(V,\omega)\to(W,\eta)\) 若满足 \(F^*\eta=\omega\)，就称为线性辛同构。标准形说明，同维数的实辛空间总是辛同构的。

\begin{corollary}\label{g4:cor:0701:top-power}
	在 \(2n\) 维空间上，交替型 \(\omega\) 非退化当且仅当 \(\omega^n\ne0\)。非退化时
	\begin{equation}\label{g4:eq:0701:symplectic-volume}
		\frac{\omega^n}{n!}
		=\mathrm{d}q^1\wedge\mathrm{d}p_1\wedge\cdots\wedge
		\mathrm{d}q^n\wedge\mathrm{d}p_n
	\end{equation}
	给出 \(V\) 的一个定向。
\end{corollary}
\begin{proof}
	在标准形中写 \(\omega=\sum_{i=1}^r e^i\wedge f^i\)。若 \(r<n\)，每个 \(n\) 次乘积都有重复的对偶向量，故 \(\omega^n=0\)。若 \(r=n\)，只有恰好各取一个二形式的乘积不消失；二形式彼此交换，故这 \(n!\) 项相同，得到式\eqref{g4:eq:0701:symplectic-volume}。
\end{proof}

上式中的正向基是交错排列的 \((e_1,f_1,\ldots,e_n,f_n)\)。若将体积形式改写成全部 \(\mathrm{d}q^i\) 在前、全部 \(\mathrm{d}p_i\) 在后的次序，就会出现因子 \((-1)^{n(n-1)/2}\)。区分坐标排列与辛定向，可以避免高维计算中的符号混淆。

\begin{example}\label{g4:exm:0701:degenerate-form}
	在 \(\mathbb R^3\) 上，\(\mathrm{d}x\wedge\mathrm{d}y\) 的根空间是 \(z\) 轴，秩为 \(2\)。在 \(\mathbb R^4\) 上，\(\mathrm{d}q^1\wedge\mathrm{d}p_1\) 的根空间由 \(\partial_{q^2},\partial_{p_2}\) 张成。后一例的环境维数为偶数，却仍然退化；必须使所有方向都进入配对，才得到辛形式。
\end{example}

\subsection{各向同性、余各向同性和 Lagrange}
\label{g4:subsec:0701:02}

固定 \(2n\) 维辛空间 \((V,\omega)\)。给定子空间 \(W\subset V\)，定义其\term[xinzhengjiaobu]{辛正交补}[symplectic orthogonal]为
\[
	W^\omega=\{v\in V:\omega(v,w)=0\text{ 对所有 }w\in W\}.
\]
它的维数与内积正交补相同，但它可能与 \(W\) 相交，甚至包含 \(W\)。

\begin{proposition}\label{g4:prop:0701:orthogonal-dimensions}
	对任意子空间 \(W\subset V\)，有
	\begin{equation}\label{g4:eq:0701:orthogonal-dimensions}
		\dim W^\omega=2n-\dim W,
		\qquad (W^\omega)^\omega=W.
	\end{equation}
	并且 \(\operatorname{rad}(\omega|_W)=W\cap W^\omega\)。因此，\(\omega|_W\) 非退化当且仅当 \(V=W\oplus W^\omega\)；此时 \(W^\omega\) 上的限制型也非退化。
\end{proposition}
\begin{proof}
	考虑线性映射 \(V\to W^*\)，它将 \(v\) 送到 \(\omega(v,\cdot)|_W\)。任一 \(W\) 上的线性泛函都可延拓为 \(V\) 上的线性泛函，而 \(\omega^\flat\) 是同构，所以该映射满射。它的核就是 \(W^\omega\)，秩与零度公式给出维数等式。

	反对称性给出 \(W\subset(W^\omega)^\omega\)，再比较维数便得双重正交补公式。限制型的根空间等于 \(W\cap W^\omega\)，这直接来自定义。它为零当且仅当 \(W,W^\omega\) 相交为零，而二者维数之和恰为 \(2n\)，故等价于直和分解。此时 \(\omega|_{W^\omega}\) 的根空间为 \(W^\omega\cap(W^\omega)^\omega=W^\omega\cap W=0\)。
\end{proof}

\begin{definition}\label{g4:def:0701:subspace-types}
	辛空间 \((V,\omega)\) 的子空间 \(W\) 称为
	\begin{enumerate}
		\item \term[gexiangtongxing]{各向同性}[isotropic]的，若 \(W\subset W^\omega\)，等价地 \(\omega|_W=0\)；
		\item \term[yugexiangtongxing]{余各向同性}[coisotropic]的，若 \(W^\omega\subset W\)；
		\item \term[lagrangezikongjian]{Lagrange 子空间}[Lagrangian subspace]\footnote{Joseph-Louis Lagrange（原名 Giuseppe Lodovico Lagrangia，1736--1813），意大利出生、后在法国工作的数学家，研究变分法、分析力学与数论。}，若 \(W=W^\omega\)；
		\item 辛子空间，若 \(\omega|_W\) 非退化。
	\end{enumerate}
\end{definition}

各向同性子空间的维数至多为 \(n\)，余各向同性子空间的维数至少为 \(n\)。正交补交换这两类子空间：\(W\) 各向同性当且仅当 \(W^\omega\) 余各向同性。Lagrange 子空间恰好同时属于这两类。

\begin{proposition}[Lagrange 子空间的判别]\label{g4:prop:0701:lagrangian-characterizations}
	对 \(W\subset V\)，下列条件等价：
	\begin{enumerate}
		\item \(W\) 是 Lagrange 子空间；
		\item \(\omega|_W=0\)，且 \(\dim W=n\)；
		\item \(W\) 是极大的各向同性子空间，即没有严格包含它的各向同性子空间。
	\end{enumerate}
	每个各向同性子空间都包含于某个 Lagrange 子空间。
\end{proposition}
\begin{proof}
	由\cref{g4:prop:0701:orthogonal-dimensions}，\(W=W^\omega\) 等价于 \(W\subset W^\omega\) 且 \(\dim W=n\)，故前两个条件等价。半维的各向同性子空间不能再扩大为各向同性子空间，因为后者的维数至多为 \(n\)。反过来，若 \(W\) 各向同性而 \(\dim W<n\)，则 \(\dim W^\omega>\dim W\)。取 \(v\in W^\omega\setminus W\)，反对称性及 \(\omega(v,v)=0\) 说明 \(W+\mathbb Rv\) 仍各向同性，所以 \(W\) 不极大。重复这一过程，有限步后便得到包含原子空间的 \(n\) 维各向同性子空间。
\end{proof}

\begin{proposition}[适配于各向同性子空间的辛基]\label{g4:prop:0701:adapted-basis}
	设 \(I\subset V\) 各向同性，\(e_1,\ldots,e_k\) 是 \(I\) 的一组基。可以将它补成一组辛基 \((e_1,\ldots,e_n,f_1,\ldots,f_n)\)。
\end{proposition}
\begin{proof}
	先将 \(I\) 扩大为 Lagrange 子空间 \(L\)，再把给定的基补成 \(L\) 的基 \(e_1,\ldots,e_n\)。由 \(V\to L^*\) 的满射性，可以选取 \(f_j\) 使 \(\omega(e_i,f_j)=\delta_{ij}\)。这些 \(f_j\) 彼此的配对还未必为零。令
	\[
		C_{ij}=\omega(f_i,f_j),\qquad
		\widehat f_j=f_j+\frac12\sum_{i=1}^n C_{ij}e_i.
	\]
	因为 \(L\) 各向同性，\(\omega(e_i,\widehat f_j)=\delta_{ij}\)；又因 \(C_{ji}=-C_{ij}\)，有
	\[
		\omega(\widehat f_i,\widehat f_j)
		=C_{ij}-\frac12C_{ij}+\frac12C_{ji}=0.
	\]
	这 \(2n\) 个向量线性无关：对它们的线性关系依次与 \(e_j\)、\(\widehat f_j\) 配对，即知所有系数为零。因此将 \(\widehat f_j\) 改记为 \(f_j\)，就得到所需辛基。
\end{proof}

\begin{example}[同为半维的两类平面]\label{g4:exm:0701:planes}
	在四维标准辛空间中，\(\operatorname{span}(e_1,e_2)\) 是 Lagrange 平面，而 \(\operatorname{span}(e_1,f_1)\) 是辛平面。两者都为半维，区别在于限制型分别为零和非退化。直线 \(\mathbb Re_1\) 各向同性，它的辛正交补 \(\operatorname{span}(e_1,e_2,f_2)\) 余各向同性；这个三维子空间上的限制型恰以 \(\mathbb Re_1\) 为根空间。
\end{example}

\begin{example}[对偶空间与对称算子的图]\label{g4:exm:0701:lagrangian-graphs}
	任意实向量空间 \(Q\) 都给出一个典范的辛空间 \(Q\oplus Q^*\)，其配对为
	\[
		\omega\bigl((q,p),(q',p')\bigr)=p'(q)-p(q').
	\]
	若一个向量与所有 \((q',0)\) 及 \((0,p')\) 的配对均为零，便依次得到 \(p=0\)、\(q=0\)，故配对非退化。\(Q\oplus0\) 和 \(0\oplus Q^*\) 都是 Lagrange 子空间。

	设 \(A:Q\to Q^*\) 线性，图空间 \(\Gamma_A=\{(q,Aq):q\in Q\}\) 已有半维。其限制配对为
	\[
		\omega\bigl((q,Aq),(q',Aq')\bigr)=(Aq')(q)-(Aq)(q').
	\]
	所以 \(\Gamma_A\) 为 Lagrange 子空间，当且仅当双线性型 \(B_A(q,q')=(Aq)(q')\) 对称。选定 \(Q\) 的基及 \(Q^*\) 的对偶基后，这等价于 \(A\) 的矩阵对称。于是，与竖直子空间 \(0\oplus Q^*\) 互补的 Lagrange 子空间，可以用全体对称矩阵来参数化。
\end{example}

上述标准形、辛正交补和图空间的描述，也可对照 Cannas da Silva 的线性辛代数综述。\cite[第1.1节，第82--84页]{Cannas2006Symplectic}

\subsection{辛约化的线性模型}
\label{g4:subsec:0701:03}

对子空间限制一个辛形式，可能产生退化方向。把这些方向商掉，就能恢复非退化性。这是辛约化的线性原型；首先说明它对任意子空间都成立。

\begin{theorem}[限制型的非退化商]\label{g4:thm:0701:linear-reduction}
	设 \(W\subset(V,\omega)\)，令 \(K=W\cap W^\omega\)，\(\pi:W\to W/K\) 为商映射。存在唯一辛形式 \(\overline\omega\) 使
	\begin{equation}\label{g4:eq:0701:reduced-form}
		\overline\omega([u],[v])=\omega(u,v),
		\qquad \pi^*\overline\omega=\omega|_W.
	\end{equation}
	特别地，若 \(C\) 余各向同性，则 \(C/C^\omega\) 是辛空间，其维数为 \(2(n-\operatorname{codim}C)\)。若 \(I\) 各向同性，则 \(I^\omega/I\) 是辛空间，其维数为 \(2(n-\dim I)\)。
\end{theorem}
\begin{proof}
	若把 \(u,v\in W\) 分别换成 \(u+k,v+\ell\)，其中 \(k,\ell\in K\)，新增的三个配对均为零。因此式\eqref{g4:eq:0701:reduced-form}不依赖代表元，并定义了交替双线性型。若 \([u]\) 与商空间中所有向量的配对为零，则 \(u\in W\) 且 \(\omega(u,v)=0\) 对所有 \(v\in W\) 成立，故 \(u\in K\)，即 \([u]=0\)。这证明非退化性。商映射满射也说明满足拉回等式的形式唯一。

	当 \(C^\omega\subset C\) 时，\(K=C^\omega\)，因此
	\[
		\dim(C/C^\omega)=\dim C-(2n-\dim C)
		=2(n-\operatorname{codim}C).
	\]
	若 \(I\subset I^\omega\)，则 \(C=I^\omega\) 余各向同性，且 \(C^\omega=I\)，代入前式即可。
\end{proof}

商空间 \((C/C^\omega,\overline\omega)\) 称为沿余各向同性子空间 \(C\) 的\term[xinyuehua]{辛约化}[symplectic reduction]。这里有两个步骤：先限制到 \(C\)，再除去限制型的根空间 \(C^\omega\)。不能直接从整个 \(V\) 商去一个非零各向同性子空间 \(I\) 并期待 \(\omega\) 自然下降，因为改变代表元 \(u\mapsto u+i\) 会使配对增加 \(\omega(i,v)\)，它对一般 \(v\in V\) 并不为零。

\ProseParagraphBreak

这个过程有明确的坐标模型。设 \(C\) 的余维为 \(k\)，则 \(I=C^\omega\) 是 \(k\) 维各向同性子空间。用\cref{g4:prop:0701:adapted-basis}选择 \(I=\operatorname{span}(e_1,\ldots,e_k)\) 的适配辛基，便有
\begin{equation}\label{g4:eq:0701:reduction-coordinates}
	C=I^\omega=\{p_1=\cdots=p_k=0\},
	\qquad
	\overline\omega=\sum_{i=k+1}^n\mathrm{d}q^i\wedge\mathrm{d}p_i.
\end{equation}
约束去掉 \(k\) 个 \(p\) 方向，取商再去掉与之配对的 \(k\) 个 \(q\) 方向。剩余的坐标成对出现，正好解释约化空间的偶数维数。

\begin{example}\label{g4:exm:0701:one-constraint}
	在 \(\mathbb R^4\) 上取 \(C=\{p_2=0\}\)。它是三维的，且 \(C^\omega=\mathbb R\partial_{q^2}\)。商映射可写为
	\[
		(q^1,q^2,p_1,0)\longmapsto(q^1,p_1),
	\]
	约化形式为 \(\mathrm{d}q^1\wedge\mathrm{d}p_1\)。在 \(C\) 内移动 \(q^2\) 不改变任何限制配对，因此这些直线应成为商空间中的单个点。若直接在整个 \(\mathbb R^4\) 中除去 \(\partial_{q^2}\)，则 \(\omega(\partial_{q^2},\partial_{p_2})=1\) 表明形式不能下降。
\end{example}

% !TeX root = ../Book4.tex
% 线性约化：在 C={p_2=0} 中沿根空间方向取商。
\begin{figure}[htbp]
	\centering
	\begin{tikzpicture}[x={(1cm,0cm)},y={(0.6cm,0.35cm)},z={(0cm,1cm)},>=Stealth,font=\small]
		\fill[GeometryBlueLight] (-1.7,-1.3,0)--(2.2,-1.3,0)--(2.2,1.7,0)--(-1.7,1.7,0)--cycle;
		\draw[GeometryGray] (-1.7,-1.3,0)--(2.2,-1.3,0)--(2.2,1.7,0)--(-1.7,1.7,0)--cycle;
		\draw[->] (-1.8,0,0)--(2.7,0,0) node[right] {\(q^1\)};
		\draw[->] (0,-1.5,0)--(0,2.1,0) node[above left] {\(p_1\)};
		\draw[->] (0,0,-0.5)--(0,0,2.8) node[above] {\(q^2\)};
		\draw[GeometryBlue,thick] (1,0.6,-0.3)--(1,0.6,2.6);
		\fill[GeometryBlue] (1,0.6,0) circle (2pt);
		\fill[GeometryBlue] (1,0.6,1.2) circle (2pt);
		\fill[GeometryBlue] (1,0.6,2.1) circle (2pt);
		\draw[GeometryBlue,->,thick] (1.24,0.6,2.1)--(1.24,0.6,0.12);
		\node[right] at (1.35,0.6,1.6) {同一个商类};
		\node[below] at (0.8,-1.2,0) {\((q^1,p_1)\) 平面};
		\node[left] at (-0.6,0,2.2) {\(C=\{p_2=0\}\)};
	\end{tikzpicture}
	\caption{三维约束空间 \(C\) 中，平行于根空间 \(\mathbb R\partial_{q^2}\) 的每条直线成为约化空间中的一个点。所画平面 \(q^2=0\) 选取了每个商类的一个代表。}
	\label{g4:fig:0701:linear-reduction}
\end{figure}


\begin{proposition}\label{g4:prop:0701:lagrangian-reduction}
	设 \(I\subset V\) 各向同性。包含 \(I\) 的 Lagrange 子空间 \(L\subset V\)，与约化辛空间 \(I^\omega/I\) 中的 Lagrange 子空间一一对应，对应为 \(L\mapsto L/I\)。
\end{proposition}
\begin{proof}
	若 \(L\) 包含 \(I\) 且各向同性，则 \(L\subset I^\omega\)。设 \(\dim I=k\)。当 \(L\) 为 Lagrange 子空间时，\(L/I\) 上的约化形式为零，且其维数为 \(n-k\)，恰为约化空间维数的一半，所以它为 Lagrange 子空间。

	反过来，若 \(\overline L\subset I^\omega/I\) 为 Lagrange 子空间，令 \(L=\pi^{-1}(\overline L)\)。拉回公式说明 \(\omega|_L=0\)，而 \(\dim L=\dim\overline L+k=n\)，所以 \(L\) 是包含 \(I\) 的 Lagrange 子空间。取像与取原像互为逆过程。
\end{proof}

\subsection{辛群及其基本拓扑}
\label{g4:subsec:0701:04}

\begin{definition}\label{g4:def:0701:symplectic-group}
	\((V,\omega)\) 的全体线性辛自同构构成\term[xinqun]{辛群}[symplectic group]，记为 \(\operatorname{Sp}(V,\omega)\)。选择辛基后，它写为
	\[
		\operatorname{Sp}(2n,\mathbb R)
		=\{S\in\operatorname{GL}(2n,\mathbb R):S^{\mathsf T}\Omega S=\Omega\}.
	\]
\end{definition}

矩阵等式直接说明，乘积和逆仍保持 \(\omega\)。以下采用矩阵空间的通常拓扑研究此群。这里的 \(\operatorname{Sp}(2n,\mathbb R)\) 是实辛群；第二卷讨论经典紧群时出现的四元数酉群 \(\operatorname{Sp}(n)\) 是另一个群，二者的拓扑不同。

\begin{proposition}\label{g4:prop:0701:symplectic-matrices}
	每个 \(S\in\operatorname{Sp}(2n,\mathbb R)\) 都有 \(\det S=1\)。若按 \((q,p)\) 分块写
	\[
		S=\begin{pmatrix}A&B\\C&D\end{pmatrix},
	\]
	则辛条件等价于
	\begin{equation}\label{g4:eq:0701:symplectic-blocks}
		A^{\mathsf T}C=C^{\mathsf T}A,\qquad
		B^{\mathsf T}D=D^{\mathsf T}B,\qquad
		A^{\mathsf T}D-C^{\mathsf T}B=I_n.
	\end{equation}
\end{proposition}
\begin{proof}
	辛变换保持 \(\omega\)，故也保持非零顶次形式 \(\omega^n/n!\)。在线性变换对顶次外幂的作用中，乘数恰为行列式，因此 \(\det S=1\)。分块等式则来自直接相乘：
	\[
		S^{\mathsf T}\Omega S=
		\begin{pmatrix}
			A^{\mathsf T}C-C^{\mathsf T}A&A^{\mathsf T}D-C^{\mathsf T}B\\
			B^{\mathsf T}C-D^{\mathsf T}A&B^{\mathsf T}D-D^{\mathsf T}B
		\end{pmatrix}.
	\]
	左下块是右上块的负转置，所以所列三个条件已包含全部要求。
\end{proof}

\begin{example}\label{g4:exm:0701:symplectic-matrices}
	在二维，任意 \(2\times2\) 矩阵都有 \(S^{\mathsf T}\Omega S=(\det S)\Omega\)，故 \(\operatorname{Sp}(2,\mathbb R)=\operatorname{SL}(2,\mathbb R)\)。在四维，按 \((q^1,q^2,p_1,p_2)\) 排列的矩阵 \(\operatorname{diag}(2,1,1,1/2)\) 虽然行列式为 \(1\)，却把标准辛形式拉回为
	\[
		2\,\mathrm{d}q^1\wedge\mathrm{d}p_1
		+\frac12\,\mathrm{d}q^2\wedge\mathrm{d}p_2,
	\]
	所以高维中的保体积条件弱于保辛条件。

	若 \(B\) 为对称矩阵，则剪切 \((q,p)\mapsto(q,p+Bq)\) 为辛变换；它将水平 Lagrange 子空间送到 \(B\) 的图。另一个辛变换族是 \((q,p)\mapsto(aq,a^{-1}p)\)，其中 \(a>0\)。当 \(a\to\infty\) 时矩阵无界，所以 \(n\ge1\) 时实辛群非紧。
\end{example}

为了把一般辛矩阵连续变为正交辛矩阵，需要对正定矩阵取实数次幂。谱定理可以逐个定义这些幂；参数依赖还需单独说明，因为不能假定重特征值附近存在光滑的特征向量基。

\begin{lemma}[正定算子的实幂及光滑依赖]\label{g4:lem:0701:positive-powers}
	设 \((E,h)\) 是有限维实内积空间，\(P:E\to E\) 自伴且正定。对每个 \(s\in\mathbb R\)，存在由谱分解确定的正定自伴算子 \(P^s\)，它在 \(P\) 的 \(\lambda\)-特征子空间上乘以 \(\lambda^s\)，并满足
	\[
		P^0=I,\qquad P^1=P,\qquad P^sP^t=P^{s+t}.
	\]
	\(P^{1/2}\) 是唯一的正定自伴平方根。任何与 \(P\) 交换的线性算子都与所有 \(P^s\) 交换；实幂与等距同构下的共轭相容。

	在正定自伴算子的开锥上，映射 \((P,s)\mapsto P^s\) 联合光滑。更一般地，若光滑向量丛上给定光滑正定内积及关于该内积自伴正定的光滑算子场，则其实幂对底点与指数联合光滑。因此此结论也适用于内积随参数变化的光滑族。
\end{lemma}
\begin{proof}
	有限维谱定理给出互相正交的特征子空间，且正定性保证所有特征值为正。在每个特征子空间上乘以 \(\lambda^s\)，就定义了不依赖基选择的 \(P^s\)。指数法则与等距共轭相容性可在各特征子空间上验证。与 \(P\) 交换的算子保持每个特征子空间，因而也与 \(P^s\) 交换。若正定自伴算子 \(B\) 满足 \(B^2=P\)，则 \(B\) 与 \(P\) 交换，且 \(B\) 在 \(P\) 的 \(\lambda\)-特征子空间上仍正定自伴，其特征值只能是 \(\sqrt\lambda\)。所以 \(B=P^{1/2}\)，证明唯一性。

	下面证明光滑性。取一组正交规范基，将算子视为对称矩阵。正定矩阵构成开集：若 \(P_0\) 的最小特征值为 \(m>0\)，则 \(\|P-P_0\|<m\) 保证 \(u^{\mathsf T}Pu>0\) 对所有非零 \(u\) 成立。固定 \(P_0\)，可选取常数 \(c>0\) 使 \(\|I-P_0/c\|<1\)，例如取 \(c\) 为 \(P_0\) 的最大特征值。在 \(P_0\) 的足够小邻域内，有统一估计 \(\|I-P/c\|\le r<1\)。于是矩阵级数
	\begin{equation}\label{g4:eq:0701:matrix-log-series}
		L(P)=(\log c)I-\sum_{k=1}^\infty\frac{(I-P/c)^k}{k}
	\end{equation}
	收敛。对矩阵变量取 \(j\) 阶导数时，乘积法则给出的项数至多为 \(k^j\)，各项在单位方向上的范数至多为 \(c^{-j}r^{k-j}\)；当 \(k<j\) 时该导数为零。因此每阶导数级数都在这个邻域内一致收敛，式\eqref{g4:eq:0701:matrix-log-series}定义光滑矩阵函数。

	逐个固定 \(P\) 再用谱定理，可以看出 \(L(P)\) 在特征值为 \(\lambda\) 的子空间上乘以 \(\log\lambda\)。所以
	\[
		P^s=\exp\bigl(sL(P)\bigr).
	\]
	矩阵指数由 \(\exp X=\sum_{k\ge0}X^k/k!\) 定义；在任意有界矩阵集上，该级数及其任意阶导数级数一致收敛，因为分母中的阶乘控制乘积求导产生的多项式因子。因此矩阵指数光滑，进而 \((P,s)\mapsto P^s\) 在 \(P_0\) 附近联合光滑。\(P_0\) 任意，得到全局结论。

	最后，在带光滑内积的向量丛上，任一局部光滑标架都可逐次正交化、再除以正长度，得到局部光滑正交规范标架。此时算子场的矩阵光滑、对称且正定，故可应用刚才的结论。在两个这样的标架之间，过渡矩阵为正交矩阵，实幂与共轭相容保证局部构造一致，从而得到全局光滑算子场。零维纤维的结论按唯一线性映射的约定成立。
\end{proof}

在标准实空间上令
\begin{equation}\label{g4:eq:0701:standard-complex-matrix}
	J_0=\begin{pmatrix}0&-I_n\\I_n&0\end{pmatrix}=-\Omega.
\end{equation}
它满足 \(J_0^2=-I\)，并且 \(\omega_0(u,J_0v)=u^{\mathsf T}v\)。将 \((q,p)\) 识别为 \(q+\mathrm{i}p\)，\(J_0\) 就是乘以 \(\mathrm{i}\)。实矩阵与 \(J_0\) 交换当且仅当它来自复线性映射；其中保持欧氏内积的矩阵恰为酉矩阵。这是三种线性结构之间最基本的联系。\cite[第3.1节，第108页]{Cannas2006Symplectic}

\begin{theorem}\label{g4:thm:0701:symplectic-retraction}
	将 \(\operatorname{U}(n)\) 通过 \(\mathbb C^n\cong\mathbb R^{2n}\) 视为实矩阵群，则
	\[
		\operatorname{Sp}(2n,\mathbb R)\cap\operatorname{O}(2n)=\operatorname{U}(n).
	\]
	而且
	\begin{equation}\label{g4:eq:0701:symplectic-retraction}
		H(S,t)=S(S^{\mathsf T}S)^{-t/2},\qquad 0\le t\le1,
	\end{equation}
	给出 \(\operatorname{Sp}(2n,\mathbb R)\) 到 \(\operatorname{U}(n)\) 的强形变收缩，即 \(H(S,0)=S\)、\(H(S,1)\in\operatorname{U}(n)\)，且对每个 \(S\in\operatorname{U}(n)\) 及每个 \(t\) 都有 \(H(S,t)=S\)。
\end{theorem}
\begin{proof}
	若 \(S\) 正交，则 \(S^{\mathsf T}\Omega S=\Omega\) 等价于 \(\Omega S=S\Omega\)，也等价于 \(J_0S=SJ_0\)。所以正交辛矩阵恰是保持欧氏内积的复线性矩阵，也就是 \(\operatorname{U}(n)\)。

	固定辛矩阵 \(S\)，令 \(P=S^{\mathsf T}S\)。首先需要证明 \(P\) 的全部实幂仍为辛矩阵。由 \(S^{\mathsf T}\Omega S=\Omega\) 取逆并整理可得
	\[
		S\Omega^{-1}S^{\mathsf T}=\Omega^{-1}.
	\]
	因为 \(\Omega^{-1}=-\Omega\)，这说明 \(S^{\mathsf T}\) 也为辛矩阵，所以 \(P=S^{\mathsf T}S\) 为正定对称辛矩阵。它满足
	\[
		P\Omega P=\Omega,
		\qquad \Omega^{-1}P\Omega=P^{-1}.
	\]
	\(\Omega\) 本身正交，故由\cref{g4:lem:0701:positive-powers}的共轭相容性，对任意实数 \(s\) 都有
	\[
		\Omega^{-1}P^s\Omega=(P^{-1})^s=P^{-s}.
	\]
	这等价于 \(P^s\Omega P^s=\Omega\)。又因 \(P^s\) 对称，所得等式正是 \(P^s\) 的辛条件。

	因此式\eqref{g4:eq:0701:symplectic-retraction}中的两个因子均为辛矩阵，\(H(S,t)\) 始终在辛群中。正定实幂的联合光滑性保证 \(H\) 连续，事实上关于矩阵和参数光滑。在 \(t=1\) 时令 \(Q=SP^{-1/2}\)，则
	\[
		Q^{\mathsf T}Q=P^{-1/2}S^{\mathsf T}SP^{-1/2}=I,
	\]
	所以 \(Q\) 属于 \(\operatorname{U}(n)\)。\(t=0\) 时为 \(S\)；若 \(S\) 原已属于 \(\operatorname{U}(n)\)，则 \(P=I\)，整个同伦都固定 \(S\)。这证明强形变收缩的全部性质。
\end{proof}

分解 \(S=QP^{1/2}\) 称为\term[jifenjie]{极分解}[polar decomposition]。其中 \(Q\) 保存欧氏长度，\(P^{1/2}\) 为正定对称算子；上述同伦把正定因子逐渐变成恒等算子。关键在于，辛条件保证正定因子的所有中间次幂也保持辛形式。

\begin{corollary}\label{g4:cor:0701:symplectic-topology}
	对 \(n\ge1\)，\(\operatorname{Sp}(2n,\mathbb R)\) 道路连通，且以恒等矩阵为基点有
	\[
		\pi_1\bigl(\operatorname{Sp}(2n,\mathbb R)\bigr)\cong\mathbb Z,
		\qquad
		\pi_2\bigl(\operatorname{Sp}(2n,\mathbb R)\bigr)=0.
	\]
	基本群的一个生成元是在最后一个 \((q^n,p_n)\) 平面上作一次正向整周旋转、在其余坐标方向上恒等的环路。
\end{corollary}
\begin{proof}
	第二卷命题8.4.2通过 \(\operatorname{U}(n-1)\to\operatorname{U}(n)\to S^{2n-1}\) 的局部平凡化与道路提升证明酉群道路连通。结合\cref{g4:thm:0701:symplectic-retraction}，任一辛矩阵先沿 \(H\) 连到酉群，再在酉群中连到单位矩阵，所以实辛群道路连通。

	强形变收缩诱导所有同伦群的同构。第二卷命题8.4.5已经由上述球面丛的长正合列证明 \(\pi_1(\operatorname{U}(n))\cong\mathbb Z\)、\(\pi_2(\operatorname{U}(n))=0\)，并证明复行列式在基本群上给出同构。其生成元为
	\[
		t\longmapsto\operatorname{diag}(1,\ldots,1,e^{2\pi\mathrm{i}t}),
		\qquad 0\le t\le1.
	\]
	把最后一个复坐标写成 \(q^n+\mathrm{i}p_n\)，这正是在最后一个实二维平面上作用为
	\[
		\begin{pmatrix}
			\cos(2\pi t)&-\sin(2\pi t)\\
			\sin(2\pi t)&\cos(2\pi t)
		\end{pmatrix}
	\]
	的环路。因此所列同伦群与生成元均由收缩得到。
\end{proof}

这个整数可以从任意辛矩阵环路 \(S(t)\) 直接读出：先取其正交因子 \(Q(t)=S(t)(S(t)^{\mathsf T}S(t))^{-1/2}\)，再计算圆周值函数 \(\det_{\mathbb C}Q(t)\) 的绕数。实行列式始终等于 \(1\)，而复行列式记录了这个环路中无法由同伦消去的整周旋转。
% 纲要条目排为编号小节标题，并进入目录与书签；说明段仍保留于纲要注释中。

% !TeX root = ../../Book4.tex
% 纲要标题：### 7.2 余切丛与典范构造
% 纲要位置：几何学计划纲要.md，第 3752 行。
% 正文状态：初稿；尚未正式审核。

% 纲要内容开始
% 
% * 典范一形式和坐标无关性
% * ω=−dθ 的符号
% * 余切提升与辛不变性
% * 闭一形式的 Lagrange 图
% 
% 纲要内容结束

\section{余切丛与典范构造}
\label{g4:sec:0702}

线性辛形式把两个方向配成一个数。要在流形上获得这样的配对，需要让它随点光滑变化，并使局部表达彼此相容。余切丛提供了一个直接的构造：余切丛的点本身就是一个线性函数，可以对投影到底流形的速度求值。由此得到的典范一形式及其外微分，把坐标变换、微分同胚和闭一形式组织到同一套几何中。这些构造也可参见 \cite[第1.3节，第85--88页；第2.1节，第93--94页]{Cannas2006Symplectic}。

\subsection{典范一形式和坐标无关性}
\label{g4:subsec:0702:01}

设 \(Q\) 为有限维、Hausdorff、第二可数的光滑无边界流形，\(\dim Q=n\)。第一卷定义11.2.9把余切空间定义为 \(\mathrm{T}_q^*Q=(\mathrm{T}_qQ)^*\)。先取这些空间的集合不交并，再用下面的坐标赋予它余切丛的光滑结构：
\[
	\mathrm{T}^*Q=\coprod_{q\in Q}\mathrm{T}_q^*Q,
	\qquad \pi_Q:\mathrm{T}^*Q\longrightarrow Q.
\]
余切丛的一点记为 \(z=(q,p)\)，其中 \(p\in\mathrm{T}_q^*Q\)。在底空间坐标 \((q^1,\ldots,q^n)\) 中，把 \(p\) 唯一写成
\[
	p=\sum_{i=1}^n p_i\,\mathrm{d}q^i|_q.
\]
这给出总空间的坐标 \((q^1,\ldots,q^n,p_1,\ldots,p_n)\)。若换用底坐标 \(Q^a=Q^a(q)\)，并写 \(p=\sum_aP_a\,\mathrm{d}Q^a|_q\)，则
\begin{equation}
	\label{g4:eq:0702:cotangent-transition}
	p_i=\sum_{a=1}^nP_a\frac{\partial Q^a}{\partial q^i},
	\qquad
	P_a=\sum_{i=1}^np_i\frac{\partial q^i}{\partial Q^a}.
\end{equation}
式\eqref{g4:eq:0702:cotangent-transition}正是第一卷命题11.2.10中余切拉回的转置变换律。过渡函数光滑，逆变换由反向换坐标给出，链式法则保证三张图上的相容性，所以这些坐标构成 \(\mathrm{T}^*Q\) 的光滑图册，其维数为 \(2n\)。不同底点上方的点由底空间的分离邻域分开，同一纤维上的不同点则由同一余切坐标图分开，故总空间 Hausdorff；从 \(Q\) 的可数坐标覆盖及各余切坐标域的可数基，也得到总空间的第二可数性。

\ProseParagraphBreak

现在取总空间的切向量 \(\Xi\in\mathrm{T}_z(\mathrm{T}^*Q)\)。它可以同时改变底点和协变向量，而 \(\mathrm{d}\pi_Q|_z(\Xi)\in\mathrm{T}_qQ\) 只保留底点的速度。恰好可以让当前的协变向量 \(p\) 对这个速度求值。

\begin{definition}
	\label{g4:def:0702:tautological-form}
	设 \(Q\) 为光滑无边界流形，\(\pi_Q:\mathrm{T}^*Q\to Q\) 为余切丛投影。其\term[dianfanyixingshi]{典范一形式}[tautological 1-form]是由
	\begin{equation}
		\label{g4:eq:0702:tautological-evaluation}
		(\theta_Q)_{(q,p)}(\Xi)
		=p\bigl(\mathrm{d}\pi_Q|_{(q,p)}\Xi\bigr),
		\qquad \Xi\in\mathrm{T}_{(q,p)}(\mathrm{T}^*Q),
	\end{equation}
	定义的一形式。
\end{definition}

下面的交换图表示的正是式\eqref{g4:eq:0702:tautological-evaluation}的求值过程；右侧箭头 \(p\) 是当前余切丛点所携带的线性函数：
\[
	\begin{tikzcd}[column sep=large,row sep=large]
		\mathrm{T}_{(q,p)}(\mathrm{T}^*Q)
		\arrow[r,"\mathrm{d}\pi_Q|_{(q,p)}"]
		\arrow[dr,swap,"(\theta_Q)_{(q,p)}"]
		& \mathrm{T}_qQ \arrow[d,"p"]\\
		& \mathbb R .
	\end{tikzcd}
\]
虽然 \(\theta_Q\) 的定义是在每一点给出的，还须核对它光滑。将切向量写成
\[
	\Xi=\sum_i a^i\frac{\partial}{\partial q^i}
	+\sum_i b_i\frac{\partial}{\partial p_i},
\]
则 \(\mathrm{d}\pi_Q(\Xi)=\sum_i a^i\partial_{q^i}\)，从而
\begin{equation}
	\label{g4:eq:0702:tautological-coordinates}
	\theta_Q=\sum_i p_i\,\mathrm{d}q^i.
\end{equation}
式\eqref{g4:eq:0702:tautological-coordinates}中的 \(\mathrm{d}q^i\) 是底坐标函数 \(q^i\circ\pi_Q\) 在总空间上的微分。这一坐标表达有光滑系数，证明了光滑性。它也直接检验了换图相容：
\[
	\sum_aP_a\,\mathrm{d}Q^a
	=\sum_i\left(\sum_aP_a\frac{\partial Q^a}{\partial q^i}\right)\mathrm{d}q^i
	=\sum_i p_i\,\mathrm{d}q^i.
\]
因此“典范”在这里有具体含义：构造只用了投影及余切向量的求值，没有选定底流形上的度量或联络。

\subsection{\texorpdfstring{\(\omega=-\mathrm{d}\theta\) 的符号}{\GeometryPdfSmallOmega{}=−d\GeometryPdfTheta{} 的符号}}
\label{g4:subsec:0702:02}

先把上一节的逐点线性条件放到流形上。以下的非退化性要求在每一点成立，闭性则涉及形式随点变化的导数。

\begin{definition}
	\label{g4:def:0702:symplectic-manifold}
	设 \(M\) 为光滑无边界流形。若光滑二形式 \(\omega\in\Omega^2(M)\) 满足 \(\mathrm{d}\omega=0\)，且在每个 \(x\in M\) 上，由
	\[
		\omega_x(v,w)=0\quad\text{对所有 }w\in\mathrm{T}_xM
	\]
	能够推出 \(v=0\)，则称 \(\omega\) 为\term[xinxingshi]{辛形式}[symplectic form]，称 \((M,\omega)\) 为\term[xinliuxing]{辛流形}[symplectic manifold]。

	给定两个辛流形 \((M,\omega_M)\)、\((N,\omega_N)\)，若光滑映射 \(F:M\to N\) 满足 \(F^*\omega_N=\omega_M\)，则称 \(F\) 为\term[xinyingshe]{辛映射}[symplectic map]；若它还是微分同胚，则称为\term[xintongpei]{辛同胚}[symplectomorphism]。
\end{definition}

上一节的交替型标准形说明辛流形的维数必为偶数。辛映射的微分一定单射：若 \(\mathrm{d}F_xv=0\)，则
\[
	\omega_M(v,w)=\omega_N(\mathrm{d}F_xv,\mathrm{d}F_xw)=0
	\quad\text{对所有 }w,
\]
故 \(v=0\)。在等维情形下，它因而是局部微分同胚；全局辛同胚的定义还包含全局可逆性。复合保持辛形式，因为 \((GF)^*=F^*G^*\)；若 \(F\) 为辛同胚，对 \(F^*\omega_N=\omega_M\) 再用 \((F^{-1})^*\)，也得到其逆保持辛形式。

\begin{theorem}
	\label{g4:thm:0702:cotangent-symplectic}
	设 \(Q\) 为光滑无边界 \(n\) 维流形。余切丛上的二形式
	\begin{equation}
		\label{g4:eq:0702:canonical-symplectic}
		\omega_Q\coloneqq-\mathrm{d}\theta_Q
		=\sum_{i=1}^n\mathrm{d}q^i\wedge\mathrm{d}p_i
	\end{equation}
	是一个辛形式，称为余切丛的\term[dianfanxinxingshi]{典范辛形式}[canonical symplectic form]。
\end{theorem}
\begin{proof}
	\(\theta_Q\) 已是全局光滑一形式。第一卷定理14.1.2给出的外微分运算与坐标无关，且 \(\mathrm{d}^2=0\)，所以 \(\omega_Q\) 全局光滑并满足 \(\mathrm{d}\omega_Q=0\)。在余切坐标中，
	\[
		-\mathrm{d}\left(\sum_i p_i\,\mathrm{d}q^i\right)
		=-\sum_i\mathrm{d}p_i\wedge\mathrm{d}q^i
		=\sum_i\mathrm{d}q^i\wedge\mathrm{d}p_i.
	\]
	若 \(\Xi=\sum_i a^i\partial_{q^i}+\sum_i b_i\partial_{p_i}\)，则
	\[
		\iota_\Xi\omega_Q=\sum_i\bigl(a^i\,\mathrm{d}p_i-b_i\,\mathrm{d}q^i\bigr).
	\]
	坐标一形式线性无关，因此这个收缩为零当且仅当全部 \(a^i,b_i\) 为零，得到非退化性。
\end{proof}

负号固定了底方向与纤维方向的配对次序：
\[
	\omega_Q(\partial_{q^i},\partial_{p_j})=\delta^i_j.
\]
这与上一节使用的矩阵
\(\begin{psmallmatrix}0&I\\-I&0\end{psmallmatrix}\)
一致。沿用 \(\theta_Q=\sum_i p_i\,\mathrm{d}q^i\) 时，取 \(\mathrm{d}\theta_Q\) 会把这个配对整体反号；本书统一采用式\eqref{g4:eq:0702:canonical-symplectic}。

\begin{example}
	\label{g4:exm:0702:fiber-reflection}
	余切反射 \(R(q,p)=(q,-p)\) 满足 \(R^*\theta_Q=-\theta_Q\)，因而 \(R^*\omega_Q=-\omega_Q\)。这种使辛形式反号的微分同胚称为\term[fanxintongpei]{反辛同胚}[anti-symplectomorphism]。当 \(\dim Q>0\) 时，\(R\) 是反辛的，不能同时保持 \(\omega_Q\)。例如在 \(\mathrm{T}^*\mathbb R\) 中，它把 \((q,p)\) 平面的有向面积反号。
\end{example}

\subsection{余切提升与辛不变性}
\label{g4:subsec:0702:03}

设 \(f:Q\to Q'\) 为微分同胚。为了把 \((q,p)\) 送到 \(f(q)\) 上方，须把 \(p\) 搬成 \(f(q)\) 处的协变向量。第一卷命题11.2.10说明，微分的对偶本来沿反方向拉回；现在利用 \(f\) 的可逆性，定义其\term[yuqietisheng]{余切提升}[cotangent lift]
\begin{equation}
	\label{g4:eq:0702:cotangent-lift}
	\widehat f:\mathrm{T}^*Q\longrightarrow\mathrm{T}^*Q',
	\qquad
	\widehat f(q,p)
	=\bigl(f(q),(\mathrm{d}f_q^{-1})^*p\bigr).
\end{equation}
这里 \(\mathrm{d}f_q^{-1}=(\mathrm{d}f_q)^{-1}
=\mathrm{d}(f^{-1})_{f(q)}\)；它从 \(\mathrm{T}_{f(q)}Q'\) 到 \(\mathrm{T}_qQ\)。等价地，像协变向量 \(p'\) 由
\[
	p'(\mathrm{d}f_qv)=p(v)\qquad(v\in\mathrm{T}_qQ)
\]
唯一确定。底映射与提升之间满足交换关系
\[
	\begin{tikzcd}[column sep=large,row sep=large]
		\mathrm{T}^*Q \arrow[r,"\widehat f"] \arrow[d,swap,"\pi_Q"]
		& \mathrm{T}^*Q' \arrow[d,"\pi_{Q'}"]\\
		Q \arrow[r,"f"] & Q'.
	\end{tikzcd}
\]

\begin{proposition}
	\label{g4:prop:0702:cotangent-lift}
	设 \(Q,Q',Q''\) 为光滑无边界流形，\(f:Q\to Q'\)、\(g:Q'\to Q''\) 为微分同胚。余切提升是光滑微分同胚，并满足
	\[
		\widehat{\operatorname{id}_Q}=\operatorname{id}_{\mathrm{T}^*Q},
		\qquad
		\widehat{g\circ f}=\widehat g\circ\widehat f,
		\qquad
		(\widehat f)^{-1}=\widehat{f^{-1}}.
	\]
	它保持典范一形式和典范辛形式：
	\[
		\widehat f^{\,*}\theta_{Q'}=\theta_Q,
		\qquad
		\widehat f^{\,*}\omega_{Q'}=\omega_Q.
	\]
\end{proposition}
\begin{proof}
	局部令 \(A(q)\) 为 \(f\) 的 Jacobian 矩阵，协变向量采用列分量，则提升的坐标式为 \((q,p)\mapsto(f(q),A(q)^{-\mathsf T}p)\)，故它光滑。恒等映射的结论直接成立。对复合，用
	\[
		\bigl(\mathrm{d}(g\circ f)_q\bigr)^{-1}
		=(\mathrm{d}f_q)^{-1}\circ(\mathrm{d}g_{f(q)})^{-1}
	\]
	并取对偶，便得到 \(\widehat{g\circ f}=\widehat g\circ\widehat f\)。令 \(g=f^{-1}\) 即得逆映射公式，也证明提升是微分同胚。

	再取 \(z=(q,p)\)、\(\Xi\in\mathrm{T}_z(\mathrm{T}^*Q)\)，记 \(\widehat f(z)=(f(q),p')\)。由交换图的微分和 \(p'\circ\mathrm{d}f_q=p\)，有
	\[
		\begin{aligned}
			(\widehat f^{\,*}\theta_{Q'})_z(\Xi)
			&=p'\bigl(\mathrm{d}\pi_{Q'}|_{\widehat f(z)}
				\,\mathrm{d}\widehat f_z\Xi\bigr)\\
			&=p'\bigl(\mathrm{d}f_q\,\mathrm{d}\pi_Q|_z\Xi\bigr)\\
			&=p\bigl(\mathrm{d}\pi_Q|_z\Xi\bigr)
			=(\theta_Q)_z(\Xi).
		\end{aligned}
	\]
	最后，第一卷命题14.1.8说明外微分与任意光滑映射的拉回交换。因此
	\[
		\widehat f^{\,*}\omega_{Q'}
		=-\widehat f^{\,*}\mathrm{d}\theta_{Q'}
		=-\mathrm{d}(\widehat f^{\,*}\theta_{Q'})
		=-\mathrm{d}\theta_Q=\omega_Q.
	\]
\end{proof}

式\eqref{g4:eq:0702:cotangent-lift}还表明，对于固定的 \(Q\)，上述函子性使 \(f\mapsto\widehat f\) 成为微分同胚群到余切丛辛同胚群的同态。它还是单射，因为 \(\pi_Q\widehat f=f\pi_Q\) 会恢复 \(f\)。这里能沿底映射的同方向提升，关键在于 \(\mathrm{d}f_q\) 可逆；对一般光滑映射，余切向量的自然运算仍是反方向的拉回。

\begin{example}
	\label{g4:exm:0702:exponential-lift}
	取 \(f:\mathbb R\to(0,\infty)\)、\(f(q)=e^q\)。在两边余切丛分别用坐标 \((q,p)\)、\((Q,P)\)，则
	\[
		\widehat f(q,p)=(e^q,e^{-q}p).
	\]
	底坐标伸长 \(e^q\) 倍，协变向量分量就缩小相同倍数，因而
	\[
		P\,\mathrm{d}Q=e^{-q}p\,e^q\mathrm{d}q=p\,\mathrm{d}q.
	\]
	进一步计算 \(\mathrm{d}Q\wedge\mathrm{d}P
	=e^q\mathrm{d}q\wedge(e^{-q}\mathrm{d}p-e^{-q}p\,\mathrm{d}q)
	=\mathrm{d}q\wedge\mathrm{d}p\)。其中由 \(p\) 的位置依赖产生的附加项，因 \(\mathrm{d}q\wedge\mathrm{d}q=0\) 消失。
\end{example}

\subsection{闭一形式的 Lagrange 图}
\label{g4:subsec:0702:04}

上一节定义了辛向量空间中的 Lagrange 子空间。对辛流形逐点使用这个定义，就得到相应的子流形概念。

\begin{definition}
	\label{g4:def:0702:lagrangian-submanifold}
	设 \((M^{2n},\omega)\) 为辛流形，\(i:L\hookrightarrow M\) 为光滑嵌入子流形。若每个 \(x\in L\) 的 \(\mathrm{T}_xL\) 都是 \((\mathrm{T}_xM,\omega_x)\) 的 Lagrange 子空间，则称 \(L\) 为\term[Lagrangeziliuxing]{Lagrange 子流形}[Lagrangian submanifold]。由\cref{g4:prop:0701:lagrangian-characterizations}，这等价于
	\[
		\dim L=n,\qquad i^*\omega=0.
	\]
\end{definition}

一个光滑一形式 \(\alpha\in\Omega^1(Q)\) 同时也是余切丛的截面。为区别它作为形式与映射的用法，记
\[
	s_\alpha:Q\longrightarrow\mathrm{T}^*Q,\qquad
	s_\alpha(q)=(q,\alpha_q),\qquad
	L_\alpha=s_\alpha(Q).
\]
等式 \(\pi_Qs_\alpha=\operatorname{id}_Q\) 说明 \(\mathrm{d}s_\alpha\) 单射，也说明 \(s_\alpha\) 到其像的逆是连续的 \(\pi_Q|_{L_\alpha}\)；所以 \(s_\alpha\) 是嵌入，\(L_\alpha\) 是一个 \(n\) 维子流形。

\begin{proposition}
	\label{g4:prop:0702:lagrangian-graph}
	设 \(Q\) 为光滑无边界 \(n\) 维流形，\(\alpha\in\Omega^1(Q)\)。则
	\begin{equation}
		\label{g4:eq:0702:section-pullback}
		s_\alpha^*\theta_Q=\alpha,\qquad
		s_\alpha^*\omega_Q=-\mathrm{d}\alpha.
	\end{equation}
	特别地，\(L_\alpha\subset(\mathrm{T}^*Q,\omega_Q)\) 为 Lagrange 子流形当且仅当 \(\alpha\) 闭。
\end{proposition}
\begin{proof}
	对 \(v\in\mathrm{T}_qQ\)，由典范一形式的定义，
	\[
		(s_\alpha^*\theta_Q)_q(v)
		=\alpha_q\bigl(\mathrm{d}\pi_Q|_{s_\alpha(q)}
			\,\mathrm{d}s_\alpha|_qv\bigr)
		=\alpha_q(v).
	\]
	外微分与拉回交换，给出第二式。图的维数已经等于 \(n\)，故其 Lagrange 条件恰为 \(s_\alpha^*\omega_Q=0\)，也就是 \(\mathrm{d}\alpha=0\)。
\end{proof}

坐标中的判别同样清楚。写 \(\alpha=\sum_i a_i(q)\,\mathrm{d}q^i\)，图由 \(p_i=a_i(q)\) 给出，则
\[
	s_\alpha^*\omega_Q
	=\sum_{i<j}
	\left(\frac{\partial a_i}{\partial q^j}
	-\frac{\partial a_j}{\partial q^i}\right)
	\mathrm{d}q^i\wedge\mathrm{d}q^j.
\]
所以图为 Lagrange 当且仅当导数矩阵 \((\partial_j a_i)\) 对称。这把上一节对线性图像的对称性判别应用到了每一点的切空间。

\ProseParagraphBreak

闭一形式在每个足够小的坐标球中都有势函数。具体地，在以 \(0\) 为中心的星形坐标域，第一卷定理14.2.2的径向公式给出
\[
	S(q)=\int_0^1\sum_i a_i(tq)q^i\,\mathrm{d}t.
\]
若 \(\partial_j a_i=\partial_i a_j\)，则
\[
	\begin{aligned}
		\partial_jS(q)
		&=\int_0^1\left(a_j(tq)+t\sum_iq^i\partial_j a_i(tq)\right)\mathrm{d}t\\
		&=\int_0^1\frac{\mathrm{d}}{\mathrm{d}t}
			\bigl(t\,a_j(tq)\bigr)\,\mathrm{d}t=a_j(q).
	\end{aligned}
\]
因此 \(L_\alpha\) 局部写成 \(p_i=\partial_iS\)。这样的 \(S\) 称为该图的\term[shengchenghanshu]{生成函数}[generating function]；两个生成函数给出同一图，当且仅当它们之差局部常值。这个结论也说明，局部生成函数可以在交叠处相差常数，而这些常数未必能通过逐片调整全部消去。

\ProseParagraphBreak

若 \(L\subset\mathrm{T}^*Q\) 为 Lagrange 子流形，并且 \(\theta_Q|_L\) 是 \(L\) 上的恰当一形式，就称它相对于 \(\theta_Q\) 为\term[qiadangLagrangeziliuxing]{恰当 Lagrange 子流形}[exact Lagrangian submanifold]。因为 \(-\theta_Q\) 是 \(\omega_Q\) 的原形式，用 \(-\theta_Q\) 作此定义得到同一个恰当性条件。对截面图，式\eqref{g4:eq:0702:section-pullback}立即给出
\[
	L_\alpha\text{ 恰当}
	\quad\Longleftrightarrow\quad
	\alpha=\mathrm{d}S\text{，其中 }S\in C^\infty(Q).
\]
这里要求的是整个 \(Q\) 上的单值函数，强于每一点附近存在势函数。

\begin{example}
	\label{g4:exm:0702:zero-fiber}
	零截面是 \(\alpha=0\) 的图，所以是恰当 Lagrange 子流形。每条余切纤维 \(\mathrm{T}_{q_0}^*Q\) 也是恰当 Lagrange 子流形：其切向量被 \(\mathrm{d}\pi_Q\) 送到零，故 \(\theta_Q\) 在纤维上为零，\(\omega_Q\) 的限制也为零；纤维的维数为 \(n\)。当 \(n>0\) 时，纤维在底空间上的投影只有一个点，因而它不可能是一形式在整个 \(Q\) 上的图。
\end{example}

\begin{example}
	\label{g4:exm:0702:circle-nonexact-graph}
	写 \(S^1=\mathbb R/(2\pi\mathbb Z)\)，令 \(\eta\) 为由 \(\mathbb R\) 上的 \(\mathrm{d}t\) 下降的一形式；平移 \(t\mapsto t+2\pi\) 保持 \(\mathrm{d}t\)，所以它确实在圆周上有定义。令 \(\pi:\mathrm{T}^*S^1\to S^1\) 为投影。用 \(p\eta_q\) 表示余切向量，便得到
	\[
		\mathrm{T}^*S^1\simeq S^1\times\mathbb R,\qquad
		\theta=p\,\pi^*\eta,\qquad \omega=\pi^*\eta\wedge\mathrm{d}p.
	\]
	对任意常数 \(c\ne0\)，\(\alpha=c\eta\) 闭，其图是圆柱上的水平圆 \(p=c\)。但沿正向圆周，
	\[
		\int_{S^1}\alpha=2\pi c\ne0.
	\]
	第一卷引理14.2.5说明恰当一形式沿闭路积分为零，所以这个 Lagrange 图不恰当。在一个角坐标小区间中仍可取 \(S=ct\)；沿圆周绕一圈以后，它增加 \(2\pi c\)，故没有下降为圆周上的单值函数。
\end{example}

闭形式图还可用纤维平移产生。给定任意 \(\alpha\in\Omega^1(Q)\)，定义
\[
	T_\alpha(q,p)=(q,p+\alpha_q).
\]
其逆为 \(T_{-\alpha}\)，因此它总是微分同胚。对 \(\Xi\in\mathrm{T}_{(q,p)}(\mathrm{T}^*Q)\)，利用 \(\pi_QT_\alpha=\pi_Q\)，直接求值得到
\[
	(T_\alpha^*\theta_Q)_{(q,p)}(\Xi)
	=(p+\alpha_q)(\mathrm{d}\pi_Q\Xi)
	=(\theta_Q+\pi_Q^*\alpha)_{(q,p)}(\Xi).
\]
于是
\begin{equation}
	\label{g4:eq:0702:fiber-translation}
	T_\alpha^*\theta_Q=\theta_Q+\pi_Q^*\alpha,\qquad
	T_\alpha^*\omega_Q=\omega_Q-\pi_Q^*(\mathrm{d}\alpha).
\end{equation}
为从式\eqref{g4:eq:0702:fiber-translation}读出充要条件，注意零截面 \(s_0\) 满足 \(s_0^*\pi_Q^*=\operatorname{id}\)，所以 \(\pi_Q^*\) 在形式上单射。因此 \(T_\alpha\) 为辛同胚当且仅当 \(\alpha\) 闭；它保持 \(\theta_Q\) 当且仅当 \(\alpha=0\)。平移把零截面送到 \(L_\alpha\)，从而再次说明闭形式的图具有 Lagrange 性。圆周例中的 \(T_{c\eta}\) 也给出一个保持辛形式、却不保持典范一形式的具体辛同胚。
% 纲要条目排为编号小节标题，并进入目录与书签；说明段仍保留于纲要注释中。

% !TeX root = ../../Book4.tex
% 纲要标题：### 7.3 相容三元组
% 纲要位置：几何学计划纲要.md，第 3759 行。
% 正文状态：初稿；尚未正式审核。

% 纲要内容开始
% 
% * 由 ω 与辅助内积构造 J
% * 正定配对 g(u,v)=ω(u,Jv)
% * 相容结构空间的可缩性
% * 光滑参数族与流形上的整体构造
% 
% 纲要内容结束

\section{相容三元组}
\label{g4:sec:0703}

辛形式测量有向面积，内积测量长度，二者起初是不同的数据。在标准平面上，把一个向量旋转四分之一周后，面积配对却变成了内积：\(\omega_0(u,J_0v)=\langle u,v\rangle\)。本节在任意辛空间中构造这种联系，并说明辅助选择可以连续地消去。构造只涉及实向量空间上的正定型和线性算子，随后逐点用于流形，不需要预先给出复坐标。

\subsection{\texorpdfstring{由 \(\omega\) 与辅助内积构造 J}{由 \GeometryPdfSmallOmega{} 与辅助内积构造 J}}
\label{g4:subsec:0703:01}

先明确“四分之一周旋转”在没有坐标时的含义。在实向量空间上，连续作用两次后变号的线性算子，可以用来定义乘以 \(\mathrm{i}\)；它不要求已有一张复坐标图。

\begin{definition}\label{g4:def:0703:almost-complex}
	设 \(V\) 为有限维实向量空间。满足 \(J^2=-I\) 的实线性算子 \(J:V\to V\) 称为 \(V\) 上的复结构。规定 \((a+\mathrm{i}b)v=av+bJv\)，便使 \(V\) 成为复向量空间。

	光滑实流形 \(M\) 上满足 \(J^2=-I\) 的光滑丛自同态 \(J:\mathrm{T}M\to\mathrm{T}M\) 称为\term[jihufujiegou]{几乎复结构}[almost complex structure]。若 \((V,\omega)\) 是辛空间，复结构 \(J\) 满足
	\begin{equation}\label{g4:eq:0703:compatibility}
		\omega(Ju,Jv)=\omega(u,v),\qquad
		\omega(u,Ju)>0\quad(u\ne0),
	\end{equation}
	则称它与 \(\omega\)\term[xiangrongfujiegou]{相容}[compatible]。流形上逐点满足这些条件的 \(J\)，称为与给定非退化二形式 \(\omega\) 相容的几乎复结构。
\end{definition}

几乎复结构的定义规定了每个切空间中的复数乘法；是否能找到使它成为标准乘法的局部复坐标，是另一个问题。当前先构造满足上述代数条件的光滑场。

\ProseParagraphBreak

固定辛空间 \((V,\omega)\)，任取一个辅助内积 \(h\)。由 \(h\) 把 \(V\) 与 \(V^*\) 识别，存在唯一算子 \(A\) 满足
\begin{equation}\label{g4:eq:0703:skew-operator}
	h(Au,v)=\omega(u,v)\qquad(u,v\in V).
\end{equation}
这里 \(h\) 对称而 \(\omega\) 反对称，故 \(A\) 应是反自伴算子。它已经具有旋转的性质，但不同方向上的尺度未必相同。下面用一个正定算子消去尺度。

\begin{theorem}\label{g4:thm:0703:polar-construction}
	设 \((V,\omega)\) 是有限维实辛空间，\(h\) 是 \(V\) 上的正定内积，\(A\) 由式\eqref{g4:eq:0703:skew-operator}定义。则 \(-A^2\) 是 \(h\)-自伴正定算子，令
	\begin{equation}\label{g4:eq:0703:normalization}
		P=(-A^2)^{1/2},\qquad
		\mathcal N_\omega(h)=J=P^{-1}A,
	\end{equation}
	其中取唯一的 \(h\)-自伴正定平方根。则 \(J\) 与 \(\omega\) 相容，且
	\begin{equation}\label{g4:eq:0703:normalized-metric}
		g(u,v)=\omega(u,Jv)=h(Pu,v)
	\end{equation}
	是正定内积。\(J\) 关于 \(h\) 正交，\(A=PJ=JP\)。
\end{theorem}
\begin{proof}
	内积的非退化性使 \(A\) 存在且唯一；\(\omega\) 非退化使 \(A\) 可逆。记 \(A^*\) 为 \(h\)-伴随，由反对称性有
	\[
		h(Au,v)=\omega(u,v)=-\omega(v,u)=-h(u,Av),
	\]
	所以 \(A^*=-A\)。从而 \(-A^2=A^*A\) 自伴，并且对 \(u\ne0\)，
	\[
		h(-A^2u,u)=h(Au,Au)>0.
	\]
	\cref{g4:lem:0701:positive-powers}给出所需的正定平方根 \(P\)。因为 \(A\) 与 \(-A^2\) 交换，它保持后者的每个特征子空间；\(P\) 在这些子空间上作用为正的标量，故 \(AP=PA\)，也有 \(AP^{-1}=P^{-1}A\)。于是
	\[
		J^2=P^{-2}A^2=-I,
		\qquad J^*=A^*P^{-1}=-P^{-1}A=-J.
	\]
	两式合起来给出 \(J^*J=I\)，即 \(h(Ju,Jv)=h(u,v)\)。又因 \(A\) 与 \(J\) 交换，
	\[
		\omega(Ju,Jv)=h(AJu,Jv)=h(JAu,Jv)=h(Au,v)=\omega(u,v).
	\]
	最后，把 \(J\) 从内积的第二个变量移到第一个变量，得到
	\[
		\omega(u,Jv)=h(Au,Jv)=h(-JAu,v)=h(Pu,v),
	\]
	其中 \(-JA=-P^{-1}A^2=P\)。\(P\) 自伴正定，故这个配对对称正定，特别地 \(\omega(u,Ju)>0\)。这证明全部相容条件与所列恒等式。
\end{proof}

分解 \(A=PJ\) 把正定尺度 \(P\) 与正交旋转 \(J\) 分开，是这里使用的极分解。构造由 \((\omega,h)\) 唯一确定；它没有为 \(\omega\) 单独指定一个优先的内积。该线性构造可对照 Cannas da Silva 的命题3.1。\cite[第3.1节，第108--109页]{Cannas2006Symplectic}

\begin{example}\label{g4:exm:0703:anisotropic-metric}
	在 \(V=\mathbb R^2\) 上取 \(\omega=\mathrm{d}q\wedge\mathrm{d}p\) 和
	\[
		h=a\,\mathrm{d}q^2+b\,\mathrm{d}p^2,\qquad a,b>0.
	\]
	在有序基 \((\partial_q,\partial_p)\) 中，直接代入定义得到
	\[
		A=\begin{pmatrix}0&-a^{-1}\\b^{-1}&0\end{pmatrix},\qquad
		P=\frac1{\sqrt{ab}}I,\qquad
		J=\begin{pmatrix}0&-\sqrt{b/a}\\\sqrt{a/b}&0\end{pmatrix}.
	\]
	因此
	\[
		g=\sqrt{a/b}\,\mathrm{d}q^2+\sqrt{b/a}\,\mathrm{d}p^2.
	\]
	当 \(a=b\) 时，\(J\partial_q=\partial_p\)、\(J\partial_p=-\partial_q\)，与本书 \(\omega\) 的正号一致；此时 \(g\) 恰为标准内积，即使辅助内积 \(h\) 是它的另一正倍数也如此。\(a/b\) 改变时，所得 \(J\) 随之改变，说明辅助内积的各向异性确实影响选择。
\end{example}

\subsection{\texorpdfstring{正定配对 \(g(u,v)=\omega(u,Jv)\)}{正定配对 g(u,v)=\GeometryPdfSmallOmega{}(u,Jv)}}
\label{g4:subsec:0703:02}

刚才构造的 \(g\) 一般不同于输入的 \(h\)。\(h\) 负责进行规范化，而 \(g\) 与输出的 \(J\) 精确满足相容恒等式。下面从这条恒等式重新辨认三个对象的关系。

\begin{proposition}\label{g4:prop:0703:compatible-triple}
	设 \((V,\omega)\) 为有限维实辛空间，\(J:V\to V\) 满足 \(J^2=-I\)。则 \(J\) 与 \(\omega\) 相容，当且仅当 \(g(u,v)=\omega(u,Jv)\) 是正定对称内积。此时
	\begin{equation}\label{g4:eq:0703:triple-identities}
		\omega(u,v)=g(Ju,v),\qquad
		g(Ju,Jv)=g(u,v),\qquad
		g(Ju,v)=-g(u,Jv).
	\end{equation}
	这样的 \((\omega,J,g)\) 称为\term[xiangrongsanyuanzu]{相容三元组}[compatible triple]。
\end{proposition}
\begin{proof}
	若 \(J\) 相容，则
	\[
		g(v,u)=\omega(v,Ju)
		=\omega(Jv,J^2u)=-\omega(Jv,u)=\omega(u,Jv)=g(u,v),
	\]
	且 \(g(u,u)>0\) 对非零 \(u\) 成立。反过来，若 \(g\) 对称正定，则正性直接给出 \(\omega(u,Ju)>0\)，而将对称性用于 \(u,Jv\) 得到
	\[
		-\omega(u,v)=g(u,Jv)=g(Jv,u)=\omega(Jv,Ju)=-\omega(Ju,Jv).
	\]
	所以 \(J\) 保持 \(\omega\)。相容性随后给出 \(g(Ju,v)=\omega(Ju,Jv)=\omega(u,v)\)，再由 \(J^2=-I\) 与反对称性得到其余两式。
\end{proof}

在相容三元组中，任意两项决定第三项：给 \(\omega,J\) 时用 \(g=\omega(\cdot,J\cdot)\)；给 \(g,J\) 时用 \(\omega=g(J\cdot,\cdot)\)；给 \(\omega,g\) 时由 \(g(Ju,v)=\omega(u,v)\) 唯一求 \(J\)。最后一种恢复的前提是这两项确实属于相容三元组；对任意辅助内积求出的算子只是 \(A\)，它还须经过式\eqref{g4:eq:0703:normalization}的规范化。

\begin{proposition}\label{g4:prop:0703:orthogonal-complex}
	设 \(V\) 为有限维实向量空间，\(g\) 是正定内积，\(J^2=-I\) 且 \(g(Ju,Jv)=g(u,v)\)。则 \(\omega(u,v)=g(Ju,v)\) 是非退化交替型，并使 \((\omega,J,g)\) 成为相容三元组。任意复结构 \(J\) 都有这样的 \(J\)-不变内积：由任意正定内积 \(h\) 可取
	\[
		g(u,v)=\tfrac12\bigl(h(u,v)+h(Ju,Jv)\bigr).
	\]
\end{proposition}
\begin{proof}
	\(J\) 的正交性与 \(J^{-1}=-J\) 给出 \(g(Ju,v)=-g(u,Jv)\)，所以 \(\omega\) 反对称。若 \(\omega(u,v)=0\) 对所有 \(v\) 成立，则 \(Ju=0\)，故 \(u=0\)。并且 \(\omega(u,Jv)=g(Ju,Jv)=g(u,v)\)，从而相容。最后所给平均显然对称正定；把 \(u,v\) 同时换成 \(Ju,Jv\)，两个加数互换，证明 \(J\)-不变性。
\end{proof}

在对应的复向量空间上，\(H(u,v)=g(u,v)-\mathrm{i}\omega(u,v)\) 是第一变量线性的 Hermitian 内积。例如 \(H(Ju,v)=\mathrm{i}H(u,v)\)，共轭对称性来自 \(g\) 对称、\(\omega\) 反对称，而 \(H(u,u)=g(u,u)>0\)。这说明三元组把同一个复内积的实部 \(g\)、负虚部 \(\omega\) 和复数乘法同时记录下来。

\begin{corollary}\label{g4:cor:0703:normalization-retraction}
	设 \((V,\omega)\) 为有限维实辛空间，\(J\) 与 \(\omega\) 相容，\(g_J(u,v)=\omega(u,Jv)\)。若用 \(h=g_J\) 作为\cref{g4:thm:0703:polar-construction}的输入，则 \(A=J\)、\(P=I\)，因而
	\[
		\mathcal N_\omega(g_J)=J.
	\]
\end{corollary}
\begin{proof}
	由式\eqref{g4:eq:0703:triple-identities}，\(g_J(Ju,v)=\omega(u,v)\)。定义 \(A\) 的唯一性给出 \(A=J\)，所以 \(-A^2=I\)，其正定平方根为 \(I\)，规范化保持 \(J\) 不变。
\end{proof}

\subsection{相容结构空间的可缩性}
\label{g4:subsec:0703:03}

记 \(\mathcal J(V,\omega)\) 为全部 \(\omega\)-相容复结构组成的集合，取它在有限维算子空间 \(\operatorname{End}_{\mathbb R}(V)\) 中的子空间拓扑。存在性已解决，接下来要比较不同选择。直接平均两个算子通常不保留平方等于 \(-I\)；正定内积却组成凸集，因此可以先平均相应内积，再恢复相容结构。

\begin{example}\label{g4:exm:0703:nonconvex-structures}
	在标准辛平面上取
	\[
		J_0=\begin{pmatrix}0&-1\\1&0\end{pmatrix},\qquad
		J_1=\begin{pmatrix}0&-1/2\\2&0\end{pmatrix}.
	\]
	二者均与 \(\mathrm{d}q\wedge\mathrm{d}p\) 相容，后者可由\cref{g4:exm:0703:anisotropic-metric}中 \(a=4,b=1\) 得到。但
	\[
		\left(\frac{J_0+J_1}{2}\right)^2=-\frac98 I,
	\]
	故其中点甚至不是复结构。相应内积为 \(g_0=\mathrm{d}q^2+\mathrm{d}p^2\)、\(g_1=2\mathrm{d}q^2+\tfrac12\mathrm{d}p^2\)，它们的任意凸组合仍正定。
\end{example}

\begin{lemma}\label{g4:lem:0703:smooth-natural-normalization}
	在一个固定有限维实向量空间 \(V\) 上，映射 \((\omega,h)\mapsto\mathcal N_\omega(h)\) 对非退化交替型 \(\omega\) 和正定内积 \(h\) 光滑。若 \(F:V\to W\) 是实线性同构，则对 \(W\) 上的此类数据 \((\omega',h')\)，有
	\[
		\mathcal N_{F^*\omega'}(F^*h')
		=F^{-1}\mathcal N_{\omega'}(h')F.
	\]
\end{lemma}
\begin{proof}
	在一个固定实基中写 \(h(u,v)=u^{\mathsf T}Hv\)、\(\omega(u,v)=u^{\mathsf T}\Omega v\)。式\eqref{g4:eq:0703:skew-operator}成为 \(A^{\mathsf T}H=\Omega\)，所以
	\[
		A=-H^{-1}\Omega.
	\]
	矩阵求逆在可逆矩阵集合上光滑，因此 \(A\) 光滑依赖 \((\Omega,H)\)。\(-A^2\) 关于变化的内积 \(h\) 自伴正定，\cref{g4:lem:0701:positive-powers}的参数版本说明 \(P=(-A^2)^{1/2}\) 也光滑；再求逆相乘即得 \(J\) 的光滑性。

	对于同构 \(F\)，定义等式直接给出源上的 \(A=F^{-1}A'F\)。\(F^{-1}P'F\) 关于 \(F^*h'\) 自伴正定，且平方为 \(-A^2\)，由正定平方根的唯一性，它就是源上的 \(P\)。于是 \(P^{-1}A=F^{-1}(P'^{-1}A')F\)，证明自然性。
\end{proof}

\begin{theorem}\label{g4:thm:0703:linear-contractibility}
	设 \((V,\omega)\) 为有限维实辛空间，固定 \(J_*\in\mathcal J(V,\omega)\)，令 \(g_*(u,v)=\omega(u,J_*v)\)。则
	\begin{equation}\label{g4:eq:0703:linear-contraction}
		\mathcal H_t(J)
		=\mathcal N_\omega\bigl((1-t)g_J+t g_*\bigr),
		\qquad 0\le t\le1,
	\end{equation}
	是 \(\mathcal J(V,\omega)\) 到 \(J_*\) 的强收缩：它连续并满足
	\[
		\mathcal H_0(J)=J,\qquad
		\mathcal H_1(J)=J_*,\qquad
		\mathcal H_t(J_*)=J_*.
	\]
	特别地，\(\mathcal J(V,\omega)\) 非空且可缩。
\end{theorem}
\begin{proof}
	对 \(u\ne0\)，\((1-t)g_J(u,u)+t g_*(u,u)>0\)，所以括号内始终是正定内积，规范化的输出始终属于 \(\mathcal J(V,\omega)\)。\(J\mapsto g_J\) 为线性公式，而规范化光滑，因此 \((t,J)\mapsto\mathcal H_t(J)\) 连续。

	在 \(t=0\) 时，\cref{g4:cor:0703:normalization-retraction}给出 \(\mathcal N_\omega(g_J)=J\)；在 \(t=1\) 时同样给出 \(\mathcal N_\omega(g_*)=J_*\)。若起点已经是 \(J_*\)，则整个度量路径恒为 \(g_*\)，故输出也一直为 \(J_*\)。这同时验证了两个端点和基点固定条件。
\end{proof}

因此任意两种相容选择都可由一条相容路径连接，且整个选择空间能连续缩到指定的一种选择。文献中的线性可缩性也可由 Lagrange 补空间和内积的参数化证明；这里采用的度量平均公式将直接用于流形。\cite[命题3.2及第3.2节，第110--112页]{Cannas2006Symplectic}

\subsection{光滑参数族与流形上的整体构造}
\label{g4:subsec:0703:04}

逐点选择相容结构，未必自动产生光滑场。上面的规范化则由同一光滑公式给出，并且与换基相容；我们只需先找到整体光滑的辅助内积。

\begin{theorem}\label{g4:thm:0703:bundle-existence}
	设 \(M\) 是有限维、Hausdorff、第二可数的光滑流形，允许有边界，\(E\to M\) 是秩为 \(2n\) 的实光滑向量丛。设 \(\omega\in\Gamma^\infty(\Lambda^2E^*)\) 在每条纤维上非退化，则存在光滑丛自同态 \(J:E\to E\)，使 \(J^2=-I\)，且 \(J\) 逐纤维与 \(\omega\) 相容。具体地，对任意光滑丛度量 \(h\)，逐纤维应用 \(J=\mathcal N_\omega(h)\) 即得这样的结构，\(g=\omega(\cdot,J\cdot)\) 也是光滑丛度量。
\end{theorem}
\begin{proof}
	第一卷第12.3节定理12.3.11保证 \(E\) 存在光滑丛度量。恢复此处需要的构造：取局部标架及从属的局部有限光滑单位分解 \(\{\chi_i\}\)，在每张标架中取标准内积 \(h_i\)，令
	\[
		h=\sum_i\chi_i h_i.
	\]
	每个加权局部张量在其定义域外补零，支集条件保证光滑；局部有限性保证每点附近只有有限项。每个非零向量的平方长度是正数的非负加权和，且至少一项权重为正，所以 \(h\) 正定。

	在任意光滑局部标架中，\(\omega,h\) 是光滑矩阵函数，\cref{g4:lem:0703:smooth-natural-normalization}给出光滑的 \(J\)。在两张标架交叠处，其自然性说明两个矩阵恰按丛自同态的规则共轭，因而表示同一个全局 \(J\)。\cref{g4:thm:0703:polar-construction}逐纤维给出 \(J^2=-I\)、相容性和正定 \(g\)，而 \(g\) 的光滑性来自张量收缩的局部公式。
\end{proof}

特别地，取 \(E=\mathrm{T}M\)，每个带有非退化光滑二形式的流形都存在与之相容的几乎复结构；辛流形当然满足这个假设。这里没有使用 \(\mathrm{d}\omega=0\)：非退化性负责逐点的构造，闭性是辛形式另外满足的微分条件。

\begin{proposition}\label{g4:prop:0703:smooth-families}
	设 \(M\) 为有限维、Hausdorff、第二可数的光滑流形，\(E\to M\) 为秩 \(2n\) 的实光滑向量丛，\(B\) 为有限维光滑参数流形；\(M,B\) 均允许有边界。若 \(\omega_b,h_b\) 是关于 \((b,x)\in B\times M\) 光滑的纤维非退化交替型与正定丛度量族，则
	\[
		J_b=\mathcal N_{\omega_b}(h_b),\qquad
		g_b=\omega_b(\cdot,J_b\cdot)
	\]
	也是关于 \((b,x)\) 光滑的相容结构与丛度量族。可以取 \(h_b\) 为一个与参数无关的固定丛度量。

	若给定 \(t\in[0,1]\) 的光滑非退化族 \(\omega_t\)，以及分别与 \(\omega_0,\omega_1\) 相容的光滑结构 \(J_0,J_1\)，则存在光滑相容族 \(J_t\)，在两个端点取得这两个指定值。
\end{proposition}
\begin{proof}
	第一部分在每个参数坐标和丛标架中直接应用\cref{g4:lem:0703:smooth-natural-normalization}，自然性保证交叠处粘合。固定辅助度量的存在性已由前一定理给出，故不需要另外作随参数的选择。

	对端点问题，令 \(g_i=\omega_i(\cdot,J_i\cdot)\)，取
	\[
		h_t=(1-t)g_0+t g_1,\qquad
		J_t=\mathcal N_{\omega_t}(h_t).
	\]
	\(h_t\) 始终正定，\(\omega_t\) 按假设始终非退化，故第一部分给出光滑族。在 \(t=i\) 处，\(h_i=g_i\)，由\cref{g4:cor:0703:normalization-retraction}，输出就是指定的 \(J_i\)。
\end{proof}

如果 \(E=\mathrm{T}M\) 且每个 \(\omega_b\) 都闭，上述结果就给出了辛形式族的相容几乎复结构族。端点构造平均的是度量；非退化的形式路径 \(\omega_t\) 是已给定的数据，并没有要求两端辛形式的任意凸组合仍非退化。

\ProseParagraphBreak

最后把可缩性提升到整个光滑截面空间。记 \(\mathcal J(E,\omega)\) 为所有光滑相容 \(J\) 的集合，在光滑自同态截面空间中取 \(C^\infty\) 紧开拓扑：局部标架中的矩阵分量及其任意有限阶导数，都在每个紧集上一致比较。这一定义不要求底流形紧致。

\begin{theorem}\label{g4:thm:0703:global-contractibility}
	设 \(M\) 为有限维、Hausdorff、第二可数的光滑流形，允许有边界，\(E\to M\) 为有限偶数秩实光滑向量丛，\(\omega\) 为光滑纤维非退化交替型。则 \(\mathcal J(E,\omega)\) 在 \(C^\infty\) 紧开拓扑下非空且可缩。固定任意 \(J_*\in\mathcal J(E,\omega)\)，令 \(g_* =\omega(\cdot,J_*\cdot)\)，一个强收缩为逐点公式
	\begin{equation}\label{g4:eq:0703:global-contraction}
		\mathcal H_t(J)(x)
		=\mathcal N_{\omega_x}\bigl((1-t)g_{J,x}+t g_{*,x}\bigr).
	\end{equation}
	对每个固定的 \(J\)，它关于 \((t,x)\) 光滑；若 \(J\) 本来随另一个有限维参数光滑变化，则所得收缩也联合光滑。若 \(J\) 与 \(J_*\) 在某个集合上相等，收缩在该集合上逐点保持 \(J_*\)。
\end{theorem}
\begin{proof}
	非空性由\cref{g4:thm:0703:bundle-existence}得到。度量的凸组合正定，规范化自然且光滑，所以式\eqref{g4:eq:0703:global-contraction}给出全局光滑相容截面及所述有限维参数光滑性。\cref{g4:cor:0703:normalization-retraction}逐点验证
	\[
		\mathcal H_0(J)=J,\qquad
		\mathcal H_1(J)=J_*,\qquad
		\mathcal H_t(J_*)=J_*.
	\]
	若在某个 \(x\) 已有 \(J_x=J_{*,x}\)，则该点的度量路径恒定，输出也恒定，证明最后的相对断言。

	还须证明在截面空间中的联合连续性，逐点连续并不足够。固定一个局部标架内的紧集 \(K\)、导数阶数 \(r\)，以及一个输入截面 \(J\)。\((1-t)g_J+t g_*\) 在紧集 \(K\times[0,1]\) 上的最小特征值有正下界。因此当另一个相容截面在 \(K\) 上充分接近 \(J\) 时，相关度量矩阵仍处在正定锥内一个与退化边界有正距离的有界区域中。固定的 \(\omega_x\) 在 \(K\) 上也有一致的非退化界。

	矩阵规范化在这些数据的邻域内是光滑函数。对式\eqref{g4:eq:0703:global-contraction}逐次求至 \(r\) 阶空间导数，链式法则把每项写为规范化映射的有限阶导数乘以输入矩阵的至 \(r\) 阶导数。所需的矩阵函数导数在上述紧数据范围内一致有界且一致连续。因此输入截面的这些导数在 \(K\) 上一致接近、参数 \(t\) 接近时，输出的至 \(r\) 阶导数也一致接近。一般紧集用有限张标架覆盖并相加估计即可。这正是 \(C^\infty\) 紧开拓扑要求的联合连续性；所需正下界只在每个紧集上选择，故论证同样适用于非紧 \(M\)。
\end{proof}

取切丛后，任意两个与同一辛形式相容的几乎复结构都有光滑相容路径相连，而全部这类选择还具有统一的连续收缩。后面用相容 \(J\) 定义依赖同伦类的对象时，选择无关性正由这个结果提供。
% 纲要条目排为编号小节标题，并进入目录与书签；说明段仍保留于纲要注释中。

% !TeX root = ../../Book4.tex
% 纲要标题：### 7.4 辛流形的整体约束
% 纲要位置：几何学计划纲要.md，第 3766 行。
% 正文状态：初稿；尚未正式审核。

% 纲要内容开始
% 
% * 闭辛流形的体积和上同调
% * 乘积、曲面和余切丛
% * 相容几乎复结构与陈类
% * 不预设 J 的可积性
% 
% ---
% 
% 纲要内容结束

\section{辛流形的整体约束}
\label{g4:sec:0704}

余切丛说明许多非紧流形自然带有辛形式，相容结构的构造则把每根切空间变成复向量空间。整体问题还要考虑这些点态资料如何与积分、上同调和坐标变化相处。首先，闭辛流形的总体积会阻止辛形式成为恰当形式；其次，相容结构的可缩性使切丛的陈类成为辛结构的不变量。最后用一个显式例子说明：切空间中的复乘法即使与辛形式相容，也未必来自复坐标。

\subsection{闭辛流形的体积和上同调}
\label{g4:subsec:0704:01}

设 \((M^{2n},\omega)\) 为辛流形。逐点取辛基 \((e_1,f_1,\ldots,e_n,f_n)\)，其对偶基记为 \((e^1,f^1,\ldots,e^n,f^n)\)。则
\[
	\frac{\omega^n}{n!}=e^1\wedge f^1\wedge\cdots\wedge e^n\wedge f^n.
\]
因此 \(\omega^n/n!\) 是处处非零的最高次形式，定义一个定向，称为\term[xindingxiang]{辛定向}[symplectic orientation]；在这个方向下它是正体积形式，称为\term[xintiji]{辛体积形式}[symplectic volume form]。任意保持 \(\omega\) 的微分同胚也保持这个定向和体积。

\begin{theorem}\label{g4:thm:0704:cohomology-obstruction}
	设 \(M\) 为非空闭连通光滑流形，实维数 \(2n\)，\(n\ge1\)，且 \(\omega\) 是其辛形式。按辛定向积分，则
	\[
		\operatorname{Vol}_\omega(M)=\frac1{n!}\int_M\omega^n>0,
		\qquad [\omega]^k\ne0\quad(1\le k\le n)
	\]
	于 \(H^{2k}_{\mathrm{dR}}(M)\) 中。特别地，\(M\) 的每个偶次 Betti 数\footnote{Enrico Betti（1823--1892），意大利数学家，研究代数、拓扑及数学物理。} \(b_{2k}\)、\(0\le k\le n\)，至少为 \(1\)，且 \(\omega\) 不恰当。
\end{theorem}
\begin{proof}
	“闭”在这里指紧且无边界。正体积形式在每张正向坐标图中是严格正函数乘坐标体积元；在一个相对紧的小图上积分已经严格为正，紧性又保证总积分有限，因此 \(\int_M\omega^n>0\)。

	因 \(\mathrm{d}\omega=0\)，每个 \(\omega^k\) 都闭，其类就是上同调乘积 \([\omega]^k\)。若某个 \(1\le k\le n\) 满足 \(\omega^k=\mathrm{d}\alpha\)，则
	\[
		\omega^n=\mathrm{d}\bigl(\alpha\wedge\omega^{n-k}\bigr).
	\]
	第一卷定理14.4.3的 Stokes 公式适用于闭定向的 \(M\)，且这里的光滑形式因紧性自动紧支，故右边在 \(M\) 上的积分为零，与正体积矛盾。于是各 \([\omega]^k\) 非零；其所在实向量空间至少一维。\(k=1\) 同时排除了 \(\omega\) 恰当的可能。
\end{proof}

这个障碍能在寻找辛形式之前使用。第二卷第7.4.2小节已由球面同调和 de Rham 积分比较算出 \(H^2_{\mathrm{dR}}(S^{2n})=0\)、\(n\ge2\)。因此四维及更高偶维球面都没有辛形式。二维球面则不同：它的面积形式就是辛形式。

\begin{example}\label{g4:exm:0704:boundary-exact}
	紧性与无边界在上一定理中各有作用。在 \(\mathbb R^{2n}\) 上，
	\[
		\omega_0=\sum_{j=1}^n\mathrm{d}q^j\wedge\mathrm{d}p_j
		=\mathrm{d}\left(\sum_{j=1}^nq^j\,\mathrm{d}p_j\right)
	\]
	是恰当辛形式。限制到一个闭球，仍是非退化闭形式；此时球虽然紧，却有边界。Stokes 公式把体积分变成边界积分，并不要求其为零。余切丛的 \(\omega=-\mathrm{d}\theta\) 是同类现象：正维底流形的全部余切纤维不紧，所以它与闭流形的障碍相容。
\end{example}

\subsection{乘积、曲面和余切丛}
\label{g4:subsec:0704:02}

在二维，任何二形式的外微分都因次数成为零。于是非退化二形式恰是面积形式；给定一个定向，任意正面积形式都提供该方向的辛结构。具体地，取一个黎曼度量并使用其定向面积形式即可。因此光滑曲面可定向，当且仅当它允许辛形式。

\begin{example}\label{g4:exm:0704:sphere-area}
	在单位球面 \(S^2\subset\mathbb R^3\) 上定义
	\[
		\omega_p(u,v)=\langle p,u\times v\rangle,
		\qquad u,v\in\mathrm{T}_pS^2.
	\]
	对非零切向量 \(u\)，向量 \(p\times u\) 也切于球面，且 \(\omega_p(u,p\times u)=|u|^2>0\)，所以 \(\omega\) 非退化。它正是外向定向的标准面积形式。在去掉两极的坐标 \((\varphi,z)\) 中，可直接求导得 \(\omega=\mathrm{d}\varphi\wedge\mathrm{d}z\)，其中 \(0\le\varphi<2\pi\)、\(-1<z<1\)。两极处的光滑性由原来的叉积公式保证，故
	\[
		\int_{S^2}\omega=\int_0^{2\pi}\int_{-1}^1\mathrm{d}\varphi\,\mathrm{d}z=4\pi.
	\]
\end{example}

\begin{proposition}\label{g4:prop:0704:products}
	设 \((M^{2m},\omega_M)\)、\((N^{2n},\omega_N)\) 是辛流形，\(\pi_M,\pi_N\) 为乘积投影。则
	\[
		\omega=\pi_M^*\omega_M+\pi_N^*\omega_N
	\]
	是 \(M\times N\) 上的辛形式，且其体积形式为
	\[
		\frac{\omega^{m+n}}{(m+n)!}
		=\pi_M^*\frac{\omega_M^m}{m!}\wedge
		\pi_N^*\frac{\omega_N^n}{n!}.
	\]
	若两因子均闭，则辛总体积为两因子辛总体积的乘积。
\end{proposition}
\begin{proof}
	拉回与外微分交换，所以 \(\omega\) 闭。在切空间分解 \(\mathrm{T}_{(x,y)}(M\times N)=\mathrm{T}_xM\oplus\mathrm{T}_yN\) 下，两因子分别配对而没有混合项。若 \((u,v)\) 与全部切向量配对为零，先用 \((u',0)\) 检验得到 \(u=0\)，再用 \((0,v')\) 检验得到 \(v=0\)，故非退化。

	二项式展开中，第一因子至多贡献 \(m\) 个二形式，第二因子至多贡献 \(n\) 个二形式，故只有包含 \(m\) 个第一项及 \(n\) 个第二项的项非零。二形式互换没有符号，留下的系数为 \(\binom{m+n}{m}\)，除以 \((m+n)!\) 即得公式。乘积定向与辛定向因此相同，分片积分与 Fubini 定理给出总体积公式。
\end{proof}

在环面 \(\mathbb R^{2n}/\mathbb Z^{2n}\) 上，坐标函数本身不全局单值，但 \(\mathrm{d}q^j,\mathrm{d}p_j\) 在整数平移下不变，因此 \(\omega_0\) 下降为辛形式。取每个坐标周期为 \(1\)，一个基本区域的辛体积为 \(1\)。任何两个闭定向曲面的乘积也给出四维闭辛流形，其二形式按上面的求和公式写出。

\ProseParagraphBreak

余切丛提供另一类普遍构造：底流形 \(Q\) 可以不可定向，\(\mathrm{T}^*Q\) 仍由典范辛形式获得定向。纤维内的动量坐标在换图时与底坐标作互相抵消的变换，正是这种定向得以成立的原因。其辛结构由底空间的光滑结构决定，构造时不需要先选度量、联络或复结构。

\subsection{相容几乎复结构与陈类}
\label{g4:subsec:0704:03}

设 \(J\) 与 \(\omega\) 相容。在实切丛上定义复数作用
\[
	(a+\mathrm{i}b)v=av+bJv,
\]
就得到复秩 \(n\) 的光滑向量丛 \((\mathrm{T}M,J)\)。它的局部标架只须对这个复数作用线性，并不要求底空间已有全纯坐标。由第二卷第21.3节的陈类构造，复秩 \(n\) 丛有自然类 \(c_k\in H^{2k}(M;\mathbb Z)\)，对丛同构不变；总类在线丛分裂后是 \(\prod_j(1+c_1(L_j))\)。下面要检查：将这些类用于切丛时，是否依赖辅助选择 \(J\)。

\begin{lemma}\label{g4:lem:0704:complex-path-isomorphism}
	设 \(E\to M\) 是光滑实向量丛，\(J_t\)、\(0\le t\le1\)，是其上光滑的复结构族。则复向量丛 \((E,J_0)\) 与 \((E,J_1)\) 同构。
\end{lemma}
\begin{proof}
	要构造一族实丛自同构 \(U_t\)，使 \(J_tU_t=U_tJ_0\)。记
	\[
		B_t=\tfrac12J_t\dot J_t.
	\]
	对 \(J_t^2=-I\) 求导，得到 \(J_t\dot J_t+\dot J_tJ_t=0\)，继而
	\[
		[B_t,J_t]=\dot J_t.
	\]
	在每根纤维上求解线性常微分方程 \(\dot U_t=B_tU_t\)、\(U_0=I\)。连续系数在有限时间区间上有界，逐次积分或线性方程的基本估计给出全区间的解；局部丛标架中的参数依赖说明解随底点光滑。不同标架中的方程按共轭转换，唯一性保证局部解粘合。逆方程 \(\dot V_t=-V_tB_t\)、\(V_0=I\)，给出 \(V_tU_t=I\)，故 \(U_t\) 处处可逆。最后
	\[
		\frac{\mathrm{d}}{\mathrm{d}t}(U_t^{-1}J_tU_t)
		=U_t^{-1}\bigl(\dot J_t+J_tB_t-B_tJ_t\bigr)U_t=0.
	\]
	其初值为 \(J_0\)，所以 \(U_1\) 正是所需的复线性丛同构。
\end{proof}

\cref{g4:thm:0703:global-contractibility}已经给出任意两个相容 \(J\) 之间的光滑道路。因此上引理及陈类的同构不变性证明下面的定义不依赖选择：
\begin{definition}\label{g4:def:0704:symplectic-chern}
	辛流形 \((M^{2n},\omega)\) 的第 \(k\) 阶陈类定义为
	\[
		c_k(M,\omega)=c_k(\mathrm{T}M,J)\in H^{2k}(M;\mathbb Z),
		\qquad 0\le k\le n,
	\]
	其中 \(J\) 是任一 \(\omega\)-相容几乎复结构。
\end{definition}

\begin{proposition}\label{g4:prop:0704:chern-invariance}
	辛同胚 \(F:(M,\omega_M)\to(N,\omega_N)\) 满足
	\[
		c_k(M,\omega_M)=F^*c_k(N,\omega_N).
	\]
	若固定流形 \(M\) 上的 \(\omega_t\)、\(0\le t\le1\)，是一族光滑辛形式，则 \(c_k(M,\omega_t)\) 不随 \(t\) 改变。
\end{proposition}
\begin{proof}
	在 \(N\) 上选相容 \(J_N\)。用 \(\mathrm{d}F\) 拉回得到 \(J_M=(\mathrm{d}F)^{-1}J_N\mathrm{d}F\)，它与 \(\omega_M\) 相容，因为相应正定配对恰为 \(F^*g_N\)。微分给出复切丛同构，再由第二卷定理21.3.2的自然性得到所述等式。

	对形式族，固定一个辅助度量，\cref{g4:prop:0703:smooth-families}的参数构造给出光滑的相容族 \(J_t\)。由\cref{g4:lem:0704:complex-path-isomorphism}，各复切丛同构，故全部陈类相同。
\end{proof}

复自然定向与辛定向一致。事实上，用相容度量逐次选单位向量及其 \(J\) 像，得到正交标架 \((e_1,Je_1,\ldots,e_n,Je_n)\)；它同时是上述两种定义中的正向标架。第二卷定理21.3.6的顶陈类与 Euler 类比较因而给出
\[
	c_n(M,\omega)=e(\mathrm{T}M)
\]
于这个共同方向下。复线丛 \(K_J=\det_{\mathbb C}(\mathrm{T}M,J)^*\) 称为该几乎复结构的\termidx[dianfanxiancong]{典范线丛}，它把第二章的典范线丛推广到了光滑几乎复情形。第二卷命题21.3.5与21.3.10的行列式及对偶公式给出
\[
	c_1(K_J)=-c_1(M,\omega).
\]
这是一条光滑复线丛；在可积情形，它才同时具有由全纯顶次形式给出的全纯线丛结构。

\begin{example}\label{g4:exm:0704:chern-models}
	标准辛环面的常向量场与标准 \(J_0\) 给出全局复标架，故所有正次陈类都为零。对于闭定向曲面 \(\Sigma_g\)，相容复切丛的秩为一。由第二卷例21.4.2的曲面切丛计算，若 \(u\) 是与面积方向一致的整数二次生成元，则
	\[
		c_1(\Sigma_g,\omega)=(2-2g)u.
	\]
	因此正面积球面的第一陈类为 \(2u\)，环面为零，亏格至少二时为负倍数；选择不同正面积形式不会改变这些值。
\end{example}

\subsection{不预设 J 的可积性}
\label{g4:subsec:0704:04}

在复坐标 \(z^j=x^j+\mathrm{i}y^j\) 中，复数乘 \(\mathrm{i}\) 对应
\[
	J\partial_{x^j}=\partial_{y^j},\qquad
	J\partial_{y^j}=-\partial_{x^j}.
\]
若某个几乎复结构在每点附近都能写成这种坐标形式，便称它\term[kejifujiegou]{可积}[integrable]。坐标之间的转移微分与 \(J\) 交换，正是多变量的 Cauchy--Riemann 条件，故这些坐标给出复流形结构。可积性要求在邻域内找坐标，相容性则只对每根切空间中的配对提出条件。

\begin{lemma}\label{g4:lem:0704:integrable-brackets}
	若光滑几乎复结构 \(J\) 来自复坐标，则任意两个复向量场 \(Z,W\) 满足 \(JZ=\mathrm{i}Z\)、\(JW=\mathrm{i}W\) 时，也有 \(J[Z,W]=\mathrm{i}[Z,W]\)。这里将 \(J\) 与 Lie 括号\footnote{Marius Sophus Lie（1842--1899），挪威数学家，创立连续变换群理论，并研究其与微分方程的联系。}复线性延拓到复化切丛。
\end{lemma}
\begin{proof}
	在复坐标中，这类向量场分别写为 \(Z=\sum a_j\partial_{z^j}\)、\(W=\sum b_j\partial_{z^j}\)，系数只须光滑。坐标向量场彼此交换，故
	\[
		[Z,W]=\sum_k\left(\sum_j a_j\partial_{z^j}b_k-b_j\partial_{z^j}a_k\right)\partial_{z^k},
	\]
	仍属于 \(\mathrm{i}\) 特征方向。此条件用内在的向量场括号表述，因而不随坐标改变。
\end{proof}

\begin{example}\label{g4:exm:0704:nonintegrable-compatible}
	在 \(\mathbb R^4\) 的坐标 \((x_1,y_1,x_2,y_2)\) 中，取
	\[
		\omega=\mathrm{d}x_1\wedge\mathrm{d}y_1+\mathrm{d}x_2\wedge\mathrm{d}y_2,
		\qquad a=e^{x_2},
	\]
	并定义
	\[
		J\partial_{x_1}=a\partial_{y_1},\quad
		J\partial_{y_1}=-a^{-1}\partial_{x_1},\quad
		J\partial_{x_2}=\partial_{y_2},\quad
		J\partial_{y_2}=-\partial_{x_2}.
	\]
	直接作用两次可得 \(J^2=-I\)。配对 \(\omega(u,Jv)\) 的矩阵为 \(\operatorname{diag}(a,a^{-1},1,1)\)，对称正定，故 \(J\) 与这个闭辛形式相容。令
	\[
		Z_1=\partial_{x_1}-\mathrm{i}a\partial_{y_1},\qquad
		Z_2=\partial_{x_2}-\mathrm{i}\partial_{y_2}.
	\]
	二者均满足 \(JZ_j=\mathrm{i}Z_j\)，然而坐标括号的乘法法则给出
	\[
		[Z_1,Z_2]=\mathrm{i}a\partial_{y_1}
		=\tfrac12(\overline Z_1-Z_1).
	\]
	右边有非零的 \(-\mathrm{i}\) 特征分量，违反\cref{g4:lem:0704:integrable-brackets}。所以这个相容 \(J\) 在任何点附近都不能由复坐标产生。
\end{example}

闭性也有独立内容。例如同一 \(\mathbb R^4\) 上的二形式
\[
	\eta=\mathrm{d}x_1\wedge\mathrm{d}y_1+e^{x_1}\,\mathrm{d}x_2\wedge\mathrm{d}y_2
\]
处处非退化，并与标准、可积的 \(J_0\) 相容，但
\[
	\mathrm{d}\eta=e^{x_1}\,\mathrm{d}x_1\wedge\mathrm{d}x_2\wedge\mathrm{d}y_2\ne0.
\]
反过来，\(\mathrm{d}x_1\wedge\mathrm{d}y_1\) 在四维空间上闭，却有二维核。这些计算分别检验了闭性、非退化性及可积性。相容的辛形式与可积复结构同时存在时，所得三元组属于 Kähler 几何\footnote{Erich Kähler（1906--2000），德国数学家，研究复几何、微分形式与外微分方程。}；第15章将系统研究可积性，第17章将展开 Kähler 结构的微分性质。当前的辛线性代数、余切构造与相容结构定理都已在光滑范畴内成立。


\begin{chaptersummary}
非退化反对称形式迫使空间为偶维，并由辛基化为标准配对。辛正交补把子空间分成各向同性、余各向同性和 Lagrange 三类；除去受限形式的核，就得到维数较小的辛空间。辛群的正定极分解提供到酉群的形变收缩，使它的拓扑可以由熟悉的紧群计算。

\ProseParagraphBreak

余切丛的典范一形式来自余切向量的求值，因而不依赖坐标。它的负外微分处处非退化；底空间微分同胚的余切提升保持这两个形式。一形式的图像为 Lagrange 子流形，当且仅当该一形式闭。这是从线性子空间进入整体子流形的第一个基本判别。

\ProseParagraphBreak

从辅助内积出发的归一化 \(J=(-A^2)^{-1/2}A\) 给出相容结构。构造随点、度量及辛形式光滑变化；先对度量插值，再归一化，得到整个相容结构空间的收缩。因此相容结构总能选取，其切丛陈类与选择无关，并在辛形式的光滑形变下保持不变。

\ProseParagraphBreak

闭辛流形的正体积迫使每个 \([\omega]^k\) 非零，排除了高维偶数球面及恰当辛形式。相容性并未自动产生复坐标：显式四维例子中，复特征向量场的括号离开原来的特征子丛，阻止了可积性。非退化性控制点态配对，闭性控制二形式的微分，可积性控制复方向能否来自坐标；后续各章将分别利用这些条件。
\end{chaptersummary}

\BookExercises
\begin{exercise}\label{g4:ex:0701:graph-action}
	在 \(\mathbb R^n_q\oplus\mathbb R^n_p\) 上取标准辛形式。设 \(S\) 为实对称矩阵，\(L_S=\{(q,Sq):q\in\mathbb R^n\}\)，且 \(M=\begin{pmatrix}A&B\\C&D\end{pmatrix}\) 为辛矩阵。若 \(A+BS\) 可逆，证明 \(M(L_S)\) 是对称矩阵
	\[
		S'=(C+DS)(A+BS)^{-1}
	\]
	的图像。对标准复结构矩阵 \(J_0=\begin{pmatrix}0&-I\\I&0\end{pmatrix}\)，求相应变换的定义范围及公式。
\end{exercise}
\begin{solution}
	图像上两个向量的配对为 \(q^{\mathsf T}Sr-r^{\mathsf T}Sq=0\)，而图像维数为 \(n\)，故 \(L_S\) 为 Lagrange 子空间。其像上的向量为
	\[
		M(q,Sq)=\bigl((A+BS)q,(C+DS)q\bigr).
	\]
	第一分量可任取，因此像恰为所述 \(S'\) 的图像。辛同构保持形式与维数，所以像仍为 Lagrange 子空间。对图像配对计算，\(u^{\mathsf T}S'v-v^{\mathsf T}S'u=0\) 对所有 \(u,v\) 成立，这正是 \(S'\) 对称。取 \(M=J_0\)，有 \(A+BS=-S\)，故这张图像坐标中的定义范围为所有可逆对称 \(S\)，变换为 \(S'=-S^{-1}\)。当 \(S\) 不可逆时，像仍是 Lagrange 子空间，只是不再能写成以整个 \(q\) 空间为底的图像。
\end{solution}

\begin{exercise}\label{g4:ex:0701:successive-reduction}
	设 \((V,\omega)\) 为线性辛空间，\(I\subset J\subset V\) 为各向同性子空间。在约化空间 \(V_I=I^\omega/I\) 中证明 \(J/I\) 各向同性，并构造保持辛形式的自然同构
	\[
		(V_I)_{J/I}\cong J^\omega/J.
	\]
	解释它为何表示分两次约化与直接约化相同。
\end{exercise}
\begin{solution}
	由 \(I\subset J\subset J^\omega\subset I^\omega\)，商 \(J/I\) 定义良好，且受限形式为零。对 \(v\in I^\omega\)，其商类与全部 \(J/I\) 配对为零，当且仅当 \(\omega(v,J)=0\)，即 \(v\in J^\omega\)。故
	\[
		(J/I)^{\overline\omega}=J^\omega/I,
		\qquad (V_I)_{J/I}=(J^\omega/I)/(J/I).
	\]
	映射 \((v+I)+(J/I)\mapsto v+J\) 不依赖代表元，满射且核为零，因而是线性同构。两边的二形式都由 \((v,w)\mapsto\omega(v,w)\) 定义，所以它保持形式。先除去 \(I\) 方向，再除去剩下的 \(J/I\) 方向，正好除去全部 \(J\) 方向；正交条件也相应从 \(I^\omega\) 收紧到 \(J^\omega\)。
\end{solution}

\begin{exercise}\label{g4:ex:0702:conformal-factor}
	设 \((M^{2n},\omega)\) 连通且辛，\(f\) 为处处非零的光滑实函数。当 \(n\ge2\) 时，证明 \(f\omega\) 为辛形式当且仅当 \(f\) 为常数；当 \(n=1\) 时说明结论如何改变。
\end{exercise}
\begin{solution}
	非零倍数保持非退化性，所以只须判断闭性。因 \(\mathrm{d}\omega=0\)，有 \(\mathrm{d}(f\omega)=\mathrm{d}f\wedge\omega\)。当 \(n\ge2\) 且该式为零时，再乘 \(\omega^{n-2}\) 得 \(\mathrm{d}f\wedge\omega^{n-1}=0\)。逐点由非退化性取唯一向量 \(v\)，使 \(\iota_v\omega=\mathrm{d}f\)。内乘的乘法法则给出
	\[
		\iota_v(\omega^n)=n\,\mathrm{d}f\wedge\omega^{n-1}=0.
	\]
	非零最高次交替形式的内乘映射是单射：若 \(v\ne0\)，将它补成一组基，体积形式在该基上的值非零。因此 \(v=0\)，进而 \(\mathrm{d}f=0\)。流形连通，\(f\) 为常数。反向显然成立。二维时一切三形式都为零，所以任何处处非零的 \(f\) 都给出辛形式；闭性此时不限制其变化。
\end{solution}

\begin{exercise}\label{g4:ex:0702:magnetic-form}
	设 \(\pi:\mathrm{T}^*Q\to Q\) 为余切丛，\(\theta\) 为典范一形式，\(\beta\) 为 \(Q\) 上的光滑二形式。证明 \(\Omega_\beta=-\mathrm{d}\theta+\pi^*\beta\) 总是非退化，并且它闭当且仅当 \(\beta\) 闭。对一形式 \(\alpha\)，计算纤维平移 \(T_\alpha(q,p)=(q,p+\alpha_q)\) 对 \(\Omega_\beta\) 的拉回；在 \(\beta\) 闭时，判定哪些一形式图像为 \(\Omega_\beta\)-Lagrange 子流形，并判断这类图像何时存在。
\end{exercise}
\begin{solution}
	在局部余切坐标中，若 \((u,v)\) 与所有向量配对为零，先用纯纤维向量 \((0,w)\) 检验，\(\Omega_\beta((u,v),(0,w))=u\cdot w\)，故 \(u=0\)。再用 \((w,0)\) 检验得到 \(-v\cdot w=0\)，所以 \(v=0\)。这证明非退化性，与 \(\beta\) 是否闭无关。外微分为 \(\mathrm{d}\Omega_\beta=\pi^*\mathrm{d}\beta\)。零截面 \(s_0\) 满足 \(s_0^*\pi^*=\mathrm{id}\)，故拉回 \(\pi^*\) 单射，得到闭性的等价条件。

	由 \(T_\alpha^*\theta=\theta+\pi^*\alpha\) 及 \(\pi T_\alpha=\pi\)，可得
	\[
		T_\alpha^*\Omega_\beta
		=-\mathrm{d}\theta+\pi^*(\beta-\mathrm{d}\alpha)
		=\Omega_{\beta-\mathrm{d}\alpha}.
	\]
	对一形式 \(\gamma\) 的图像截面，\(s_\gamma^*\Omega_\beta=\beta-\mathrm{d}\gamma\)。图像维数为总空间的一半，故它为 Lagrange 子流形当且仅当 \(\mathrm{d}\gamma=\beta\)。所以这种全局图像存在当且仅当 \(\beta\) 恰当。局部闭形式总有原函数一形式，故局部图像与全局图像的存在条件不同。
\end{solution}

\begin{exercise}\label{g4:ex:0702:conormal}
	设 \(S^k\subset Q^n\) 是嵌入子流形。定义其\term[yufacong]{余法丛}[conormal bundle]为
	\[
		N^*S=\{(q,p)\in\mathrm{T}^*Q:q\in S,\ p|_{\mathrm{T}_qS}=0\}.
	\]
	证明 \(N^*S\) 为 \(\mathrm{T}^*Q\) 的 Lagrange 子流形。将结论用于 \(Q\times Q\) 中的对角线，写出所得余法丛。
\end{exercise}
\begin{solution}
	在适配坐标中，\(S\) 由 \(q^{k+1}=\cdots=q^n=0\) 描述，其余法条件为 \(p_1=\cdots=p_k=0\)。因此 \(N^*S\) 是光滑嵌入子流形，自由坐标为 \((q^1,\ldots,q^k,p_{k+1},\ldots,p_n)\)，维数恰为 \(n\)。记包含映射为 \(j\)。对 \(\xi\in\mathrm{T}_{(q,p)}N^*S\)，\(\mathrm{d}\pi(\xi)\in\mathrm{T}_qS\)，所以
	\[
		(j^*\theta)(\xi)=p\bigl(\mathrm{d}\pi(\xi)\bigr)=0.
	\]
	故 \(j^*\omega=-\mathrm{d}(j^*\theta)=0\)，结合维数即得 Lagrange 条件。对角线的切向量为 \((u,u)\)，协向量 \((p,r)\) 在它们上恒为零等价于 \(p+r=0\)。于是所求余法丛为 \(\{(q,q,p,-p):q\in Q,\ p\in\mathrm{T}_q^*Q\}\)。
\end{solution}

\begin{exercise}\label{g4:ex:0703:plane-parameters}
	在 \(\mathbb R^2_{q,p}\) 上固定 \(\omega=\mathrm{d}q\wedge\mathrm{d}p\)。证明相容复结构的全体与上半平面 \(\{x+\mathrm{i}y:y>0\}\) 光滑一一对应，公式为
	\[
		g_{x,y}=\frac1y\begin{pmatrix}1&x\\x&x^2+y^2\end{pmatrix},
		\qquad J_{x,y}=\frac1y\begin{pmatrix}-x&-(x^2+y^2)\\1&x\end{pmatrix}.
	\]
	求一个与 \(J_{x,y}\) 对应的线性复坐标。
\end{exercise}
\begin{solution}
	令 \(\Omega=\begin{pmatrix}0&1\\-1&0\end{pmatrix}\)。相容度量矩阵为 \(G=\Omega J\)，所以 \(J=-\Omega G\)。写正定对称矩阵为 \(G=\begin{pmatrix}a&b\\b&c\end{pmatrix}\)，直接计算得到 \((-\Omega G)^2=-(ac-b^2)I\)。因此相容性等价于 \(G\) 正定且 \(\det G=1\)。此时取 \(y=1/a>0\)、\(x=b/a\)，由 \(c=(1+b^2)/a\) 恢复题中公式，反向代入也显然正定且行列式为1。参数提取与恢复都光滑，给出所述对应。

	取 \(\tau=x+\mathrm{i}y\)、\(z=q+\tau p\)。其作为实线性映射的行列式为 \(y>0\)，因而是复坐标。矩阵相乘可核对 \(\mathrm{d}z\circ J_{x,y}=\mathrm{i}\,\mathrm{d}z\)，且
	\[
		g_{x,y}=\frac1y|\mathrm{d}q+\tau\,\mathrm{d}p|^2.
	\]
	这把同一面积形式的所有相容复结构直接写成了一族平面形状。
\end{solution}

\begin{exercise}\label{g4:ex:0703:tamed-not-compatible}
	称几乎复结构 \(J\) 被二形式 \(\omega\) \term[xunfu]{驯服}[tamed]，是指 \(\omega(v,Jv)>0\) 对每个非零实向量成立。在 \(\mathbb R^4\) 上取标准复坐标 \(z_j=x_j+\mathrm{i}y_j\)、标准 \(J_0\)，并令
	\[
		\omega_t=\mathrm{d}x_1\wedge\mathrm{d}y_1+\mathrm{d}x_2\wedge\mathrm{d}y_2
		+t\operatorname{Re}(\mathrm{d}z_1\wedge\mathrm{d}z_2),\qquad t\in\mathbb R.
	\]
	证明每个 \(\omega_t\) 都是驯服 \(J_0\) 的辛形式，但 \(t\ne0\) 时并不与 \(J_0\) 相容。
\end{exercise}
\begin{solution}
	记 \(\beta=\operatorname{Re}(\mathrm{d}z_1\wedge\mathrm{d}z_2)=\mathrm{d}x_1\wedge\mathrm{d}x_2-\mathrm{d}y_1\wedge\mathrm{d}y_2\)。由 \(\mathrm{d}z_j(J_0u)=\mathrm{i}\,\mathrm{d}z_j(u)\)，有 \(\beta(J_0u,J_0v)=-\beta(u,v)\)。令第二向量为 \(J_0v\)，利用反对称性可得 \(\beta(v,J_0v)=-\beta(v,J_0v)\)，故它为零。因此
	\[
		\omega_t(v,J_0v)=\omega_0(v,J_0v)=|v|^2.
	\]
	这证明驯服性，也排除了 \(\omega_t\) 的非零核。各系数为常数，故 \(\omega_t\) 闭，从而辛。另一方面，取 \(u=\partial_{x_1}\)、\(v=\partial_{x_2}\)，有 \(\omega_t(u,v)=t\)，却有 \(\omega_t(J_0u,J_0v)=-t\)。当 \(t\ne0\) 时不满足相容性要求的不变条件。这表明正性条件本身不能代替完整的相容条件。
\end{solution}

\begin{exercise}\label{g4:ex:0704:reverse-form}
	设 \((M^{2n},\omega)\) 为辛流形。证明
	\[
		c_k(M,-\omega)=(-1)^k c_k(M,\omega),
	\]
	并比较 \(\omega\) 与 \(-\omega\) 的辛定向。在 \(S^2\times S^2\) 上取 \(\omega=\pi_1^*\sigma+\pi_2^*\sigma\)，其中 \(\sigma\) 为正面积形式，证明 \(\omega\) 与 \(-\omega\) 虽给出相同定向，却不能由固定流形上的光滑辛形式族连接。
\end{exercise}
\begin{solution}
	若 \(J\) 与 \(\omega\) 相容，则 \(-J\) 与 \(-\omega\) 相容，相应度量仍为 \(g\)。复丛 \((\mathrm{T}M,-J)\) 是 \((\mathrm{T}M,J)\) 的共轭丛；第二卷命题21.3.10的共轭公式给出所述陈类等式。又因 \((-\omega)^n=(-1)^n\omega^n\)，二者在 \(n\) 为偶数时定向相同，为奇数时定向相反。

	在两球面的乘积上，令 \(x,y\) 为分别从两因子拉回的正整数二次生成元。乘积相容结构使复切丛分成两条拉回线丛，故由球面的 \(c_1=2u\) 和总陈类的直和公式，得到 \(c_1(M,\omega)=2x+2y\)，而 \(c_1(M,-\omega)=-2x-2y\)。两类不同，例如在 \(S^2\times\{q\}\) 上的积分分别为 \(2\) 与 \(-2\)。若有光滑辛形式族连接它们，\cref{g4:prop:0704:chern-invariance} 要求第一陈类相同，矛盾。这里实维数为4，故定向相同，不能单凭定向区别这两个辛形式。
\end{solution}

\begin{exercise}\label{g4:ex:0704:cotangent-chern}
	设 \(Q\) 为可定向光滑流形，\(\mathrm{T}^*Q\) 配备典范辛形式。证明 \(c_1(\mathrm{T}^*Q)=0\)。提示：先把切丛限制到零截面，并借助底空间度量识别得到的复向量丛。
\end{exercise}
\begin{solution}
	记零截面为 \(s\)。沿零截面，微分 \(\mathrm{d}s\) 和竖直纤维给出自然分裂
	\[
		s^*\mathrm{T}(\mathrm{T}^*Q)\cong\mathrm{T}Q\oplus\mathrm{T}^*Q,
		\qquad \omega((u,p),(v,r))=r(u)-p(v).
	\]
	选底空间度量 \(h\)，并以同一字母记 \(h:\mathrm{T}Q\to\mathrm{T}^*Q\)。定义 \(J_h(u,p)=(-h^{-1}p,hu)\)，则 \(J_h^2=-I\)，且相应度量为 \(h(u,v)+h^{-1}(p,r)\)，所以相容。复线性同构为
	\[
		(u,p)\longmapsto u+\mathrm{i}h^{-1}p
		\quad\hbox{到}\quad\mathrm{T}Q\otimes_{\mathbb R}\mathbb C.
	\]
	任一全局相容结构限制到零截面后，都与 \(J_h\) 由光滑相容族连接；\cref{g4:lem:0704:complex-path-isomorphism} 保证所得复丛同构。因此第一陈类等于复化切丛的第一陈类，即其复行列式线丛的第一陈类。该行列式线丛为 \((\det_{\mathbb R}\mathrm{T}Q)\otimes_{\mathbb R}\mathbb C\)。可定向性使实行列式线丛具有处处非零截面，从而复化后也平凡，故 \(s^*c_1(\mathrm{T}^*Q)=0\)。最后纤维缩放 \((q,p)\mapsto(q,tp)\) 给出余切丛到零截面的形变收缩，\(s^*\) 为上同调同构，所以原来的第一陈类也为零。
\end{solution}

\begin{exercise}\label{g4:ex:0704:complex-submanifold}
	设 \(J\) 与 \((M,\omega)\) 相容，嵌入子流形 \(N\subset M\) 满足 \(J(\mathrm{T}N)=\mathrm{T}N\)。证明 \(N\) 为辛子流形，并证明它的辛正交法丛也是 \(J\)-不变的。用这两个复丛写出 \(c(\mathrm{T}M)|_N\) 的乘积公式。
\end{exercise}
\begin{solution}
	若 \(0\ne u\in\mathrm{T}_xN\)，则 \(Ju\in\mathrm{T}_xN\)，且 \(\omega(u,Ju)=g(u,u)>0\)。所以限制形式非退化；它的闭性来自拉回与外微分交换，故 \(N\) 辛。设 \(E=(\mathrm{T}N)^\omega\)。因 \(\mathrm{T}N\) 的限制形式非退化，\(\mathrm{T}N\cap E=0\)，维数公式给出 \(\mathrm{T}M|_N=\mathrm{T}N\oplus E\)。这是光滑丛分解，因为 \(E\) 也等于 \(g\)-正交补：\(g(v,w)=\omega(v,Jw)\)，而 \(J\) 保持 \(\mathrm{T}N\)。若 \(v\in E\)、\(w\in\mathrm{T}N\)，则 \(\omega(Jv,w)=-\omega(v,Jw)=0\)，故 \(Jv\in E\)。于是两项均为复子丛，由第二卷定理21.3.7得
	\[
		c(\mathrm{T}M,J)|_N=c(\mathrm{T}N,J|_{\mathrm{T}N})\,c(E,J|_E).
	\]
	尤其第一陈类相加。这里仅用到切空间的 \(J\)-不变性，没有要求 \(J\) 可积。
\end{solution}
